\documentclass[11pt, a4paper, oneside]{book}

% ── Engine: XeLaTeX ──────────────────────────────────────────────────────────
\usepackage{fontspec}
\setmainfont{Latin Modern Roman}
\setsansfont{Latin Modern Sans}
\setmonofont{DejaVu Sans Mono}[Scale=0.85]

% ── Core packages ────────────────────────────────────────────────────────────
\usepackage[margin=2.5cm]{geometry}
\usepackage{amsmath, amssymb, amsthm, mathtools}
\usepackage{tikz}
\usetikzlibrary{positioning, arrows.meta, calc, matrix, fit, backgrounds,
  decorations.pathreplacing, shapes.geometric}
\usepackage[outputdir=build]{minted}
\setminted{
  fontsize=\small,
  breaklines=true,
  bgcolor=codebg,
  frame=none,
  framesep=2mm,
  baselinestretch=1.15,
}
\usepackage{xcolor}
\usepackage{hyperref}
\hypersetup{colorlinks=true, linkcolor=deepblue, urlcolor=midblue, citecolor=accentgreen}
\usepackage{microtype}
\usepackage{enumitem}
\usepackage{booktabs}
\usepackage{tcolorbox}
\tcbuselibrary{theorems, skins, breakable}
\usepackage{fancyhdr}
\usepackage{sectsty}   % heading font control WITHOUT reimplementing \part
\usepackage{epigraph}
\usepackage{array}
\usepackage{longtable}
\usepackage{tabularx}
\usepackage{multirow}

% ── Colour palette ───────────────────────────────────────────────────────────
\definecolor{deepblue}{HTML}{1a3a5c}
\definecolor{midblue}{HTML}{2e6da4}
\definecolor{accentgreen}{HTML}{2a7a4b}
\definecolor{accentorange}{HTML}{c0622a}
\definecolor{accentpurple}{HTML}{5e3c8f}
\definecolor{accentred}{HTML}{b5332e}
\definecolor{codebg}{HTML}{f5f5f0}
\definecolor{defbg}{HTML}{eaf2fb}
\definecolor{thmbg}{HTML}{eafaf1}
\definecolor{exbg}{HTML}{fef9e7}
\definecolor{histbg}{HTML}{f3eef8}
\definecolor{warnbg}{HTML}{fdecea}
\definecolor{lightblue}{HTML}{dce9f5}
\definecolor{lightgreen}{HTML}{d4efdf}
\definecolor{lightorange}{HTML}{fdebd0}
\definecolor{lightred}{HTML}{f9e0df}

% ── tcolorbox environments ───────────────────────────────────────────────────
\newtcolorbox{defbox}[1][]{
  enhanced, breakable,
  colback=defbg, colframe=midblue!70,
  coltitle=black, fonttitle=\bfseries,
  title={Definition},
  #1
}

\newtcolorbox{thmbox}[1][]{
  enhanced, breakable,
  colback=thmbg, colframe=accentgreen!70,
  coltitle=black, fonttitle=\bfseries,
  title={Theorem},
  #1
}

\newtcolorbox{lembox}[1][]{
  enhanced, breakable,
  colback=thmbg, colframe=accentgreen!50,
  coltitle=black, fonttitle=\bfseries,
  title={Lemma},
  #1
}

\newtcolorbox{propbox}[1][]{
  enhanced, breakable,
  colback=thmbg, colframe=accentgreen!50,
  coltitle=black, fonttitle=\bfseries,
  title={Proposition},
  #1
}

\newtcolorbox{histbox}[1][]{
  enhanced, breakable,
  colback=histbg, colframe=accentpurple!50,
  coltitle=black, fonttitle=\bfseries,
  title={Historical Note},
  #1
}

\newtcolorbox{exercisebox}[1][]{
  enhanced, breakable,
  colback=lightorange!50, colframe=accentorange!70,
  coltitle=black, fonttitle=\bfseries,
  title={Exercise},
  #1
}

\newtcolorbox{leanbox}[1][]{
  enhanced, breakable,
  colback=codebg, colframe=midblue!60,
  coltitle=black, fonttitle=\bfseries\ttfamily,
  title={Lean 4},
  #1
}

\newtcolorbox{intuitionbox}[1][]{
  enhanced, breakable,
  colback=lightblue!40, colframe=midblue!40,
  coltitle=black, fonttitle=\bfseries,
  title={Intuition},
  #1
}

\newtcolorbox{warningbox}[1][]{
  enhanced, breakable,
  colback=warnbg, colframe=accentred!60,
  coltitle=black, fonttitle=\bfseries,
  title={Warning},
  #1
}

\newtcolorbox{crossrefbox}[1][]{
  enhanced, breakable,
  colback=histbg!50, colframe=accentpurple!30,
  coltitle=black, fonttitle=\bfseries,
  title={Cross-Reference: Lattice Book},
  #1
}

% ── Custom macros ────────────────────────────────────────────────────────────
\newcommand{\Z}{\mathbb{Z}}
\newcommand{\Q}{\mathbb{Q}}
\newcommand{\R}{\mathbb{R}}
\newcommand{\N}{\mathbb{N}}
\newcommand{\Finset}{\operatorname{Finset}}
\newcommand{\dom}{\operatorname{dom}}
\newcommand{\obj}{\operatorname{obj}}
\newcommand{\pen}{\operatorname{pen}}
\newcommand{\val}{\operatorname{val}}
\newcommand{\statTotal}{\operatorname{statTotal}}
\newcommand{\maxContrib}{\operatorname{maxContrib}}
\newcommand{\minContrib}{\operatorname{minContrib}}
\newcommand{\opt}{\operatorname{opt}}
\newcommand{\LP}{\operatorname{LP}}
\newcommand{\IP}{\operatorname{IP}}
\DeclareMathOperator*{\argmin}{arg\,min}

% ── Chapter/section styling ──────────────────────────────────────────────────
% ── Heading colors via sectsty ────────────────────────────────────────────────
\allsectionsfont{\bfseries}

% ── Headers and footers ─────────────────────────────────────────────────────
\pagestyle{fancy}
\fancyhf{}
\fancyhead[L]{\small\itshape\leftmark}
\fancyhead[R]{\small\thepage}
\renewcommand{\headrulewidth}{0.4pt}

% ════════════════════════════════════════════════════════════════════════════
\begin{document}

% ── Title page ───────────────────────────────────────────────────────────────
\begin{titlepage}
\centering
\vspace*{3cm}
{\Huge\bfseries\color{deepblue} Verified Optimization\\[0.4em]}
{\LARGE\color{midblue} From Lattices to Solvers in Lean\,4\\[2em]}
{\large\color{accentgreen} A Zero-to-Hero Guide to Combinatorial Optimization\\
with a GW2 Build Editor Capstone\\[4em]}
\vfill
{\large\itshape You built a propagator network for gear optimization.\\
It prunes beautifully. Now let's make it \textbf{solve}.\\[3em]}
\vfill
{\small Companion to \emph{Verified Order Theory} (the lattice/propagator book)}
\end{titlepage}

% ── Table of contents ────────────────────────────────────────────────────────
\tableofcontents
\newpage

% ── Preface ──────────────────────────────────────────────────────────────────
\chapter*{Preface}
\addcontentsline{toc}{chapter}{Preface}

This book has an unusual origin.

It began with a practical problem: building a gear optimizer for Guild Wars~2's
World vs.\ World game mode.  You choose stat combinations for 16+ equipment
slots, each drawn from a finite menu of options.  You have stat targets to meet,
a rune constraint that couples all six armor pieces, and an infusion budget for
fine-tuning.  The goal: find the loadout that gets as close to your targets as
possible, penalising shortfalls much more than overshoots.

The first attempt was a \emph{propagator network}---the architecture from our
companion book \emph{Verified Order Theory: From Partial Orders to Propagator
Networks in Lean\,4}.  Each gear slot became a cell holding a set of remaining
stat combos, ordered by reverse inclusion.  Propagators enforced constraints:
the rune propagator ensured all six armor rune slots matched; the stat-floor
propagator detected when no remaining combination of gear could reach a target.
Running the network to its Knaster--Tarski fixed point pruned the search space
dramatically.

But propagation alone couldn't \emph{solve} the problem.  It could tell you
that Dire gear on the headgear slot was impossible given your power target, but
it couldn't tell you whether Marauder or Berserker on that slot led to a better
overall build.  Propagation eliminates the impossible; it doesn't choose the
optimal among the possible.

This book picks up where propagation left off.  It gives the propagator engine
a driver: systematic search to explore the pruned space, bounding techniques to
skip parts of the search that provably can't beat what you've already found, and
heuristic methods to find good starting solutions fast.  The result is a
\emph{verified solver}---a Lean\,4 program that either returns an optimal build
with a machine-checkable proof of optimality, or certifies that no feasible
build exists.

The mathematics is real: linear programming, integer programming, branch and
bound, constraint propagation, Pareto optimality, dynamic programming.  But
every concept is motivated by the gear problem, implemented in Lean\,4, and
connected back to the lattice-theoretic foundations you already know.

\subsection*{Prerequisites}

Comfort with Lean\,4 typeclasses, basic tactic proofs, and the material from
\emph{Verified Order Theory}---especially partial orders, lattices,
join-semilattices, propagator cells, and the Knaster--Tarski fixed-point
theorem.

\subsection*{How to Read This Book}

Part~I opens with the build editor and your propagator attempt, then
strengthens the propagator engine before introducing any new techniques.
Part~II develops the mathematical tools: linear programming, integer
programming, search strategies.  Part~III adds heuristic methods and
multi-objective optimization.  Part~IV assembles everything into the capstone
solver with a formal correctness theorem.

Cross-reference boxes connect each new concept to its lattice-theoretic
counterpart.  You'll find that optimisation and order theory are deeply
intertwined---more so than most textbooks suggest.

\bigskip\noindent
Let's build.

\newpage

% ════════════════════════════════════════════════════════════════════════════
%  PART I — THE PROBLEM AND THE ENGINE YOU ALREADY HAVE
% ════════════════════════════════════════════════════════════════════════════
\part{The Problem and the Engine You Already Have}

% ────────────────────────────────────────────────────────────────────────────
\chapter{The Build Editor and the Propagator Attempt}
\label{ch:problem}
% ────────────────────────────────────────────────────────────────────────────

\epigraph{The question is not whether we can prune the search space,\\
but whether we can navigate what remains.}{}

\section{The Gear Optimisation Problem}

In Guild Wars~2's World vs.\ World mode, a character's combat effectiveness is
determined by nine \emph{attributes} (stats).  Each stat has a different effect:

\begin{center}
\begin{tabular}{ll}
\toprule
\textbf{Stat} & \textbf{Effect} \\
\midrule
Power & Increases direct damage \\
Precision & Increases critical strike chance \\
Ferocity & Increases critical strike damage \\
Toughness & Increases armour \\
Vitality & Increases maximum health \\
Condition Damage & Increases damage-over-time \\
Expertise & Increases condition duration \\
Concentration & Increases boon duration \\
Healing Power & Increases outgoing healing \\
\bottomrule
\end{tabular}
\end{center}

A character's stats come from three sources: base stats (determined by class),
gear, and temporary buffs (food, utility, traits).  This book concerns the
gear component, which the player has full control over.

\subsection{Equipment Slots}

A full equipment loadout consists of:

\paragraph{Armour (6 pieces).} Headgear, shoulders, coat, gloves, leggings,
boots.  Each piece has one rune slot and one infusion slot.

\paragraph{Weapons.}  The active weapon configuration is either:
\begin{itemize}[nosep]
  \item one two-handed weapon (2 sigil slots, 2 infusion slots), or
  \item one main-hand weapon + one off-hand weapon (1 sigil and 1 infusion each).
\end{itemize}
A character may have two weapon sets, but only the active set's stats count.
We optimise only the active set.

\paragraph{Trinkets.}  Amulet ($\times 1$, 0 infusions), ring ($\times 2$,
3 infusions each), accessory ($\times 2$, 1 infusion each), back piece
($\times 1$, 2 infusions).

The total infusion budget is $6 + 2 + 6 + 2 + 2 = 18$ slots, each contributing
$+5$ to a single stat of the player's choice.  That's 90 freely distributable
stat points.

\subsection{Stat Combinations}

Each gear piece is assigned a \emph{stat combination} (also called a
\emph{prefix}).  There are 40 WvW-legal stat combinations at Ascended rarity,
falling into three categories:

\paragraph{Triple-attribute (23 combos).}  One \emph{major} stat and two
\emph{minor} stats.  Major stats receive a larger coefficient; minor stats
receive a smaller one.

\paragraph{Quadruple-attribute (16 combos).}  Two major stats and two minor
stats.  Each major stat gets the 4-stat major coefficient; each minor stat gets
the 4-stat minor coefficient.

\paragraph{Celestial (1 combo).}  All seven of Power, Precision, Ferocity,
Vitality, Toughness, Healing Power, and Condition Damage receive the Celestial
coefficient (equal for all).  The WvW version of Celestial excludes Expertise
and Concentration (nerfed relative to PvE).

\medskip

The coefficients depend on the gear slot.  At Ascended rarity and level~80:

\begin{center}
\small
\begin{tabular}{lccccc}
\toprule
\textbf{Slot} & \textbf{3-stat Maj} & \textbf{3-stat Min} &
  \textbf{4-stat Maj} & \textbf{4-stat Min} & \textbf{Celestial} \\
\midrule
\multicolumn{6}{l}{\textit{Armour}} \\
\quad Headgear     & 63  & 45  & 54  & 30 & 30 \\
\quad Shoulders    & 47  & 34  & 40  & 22 & 22 \\
\quad Coat         & 141 & 101 & 121 & 67 & 67 \\
\quad Gloves       & 47  & 34  & 40  & 22 & 22 \\
\quad Leggings     & 94  & 67  & 81  & 44 & 44 \\
\quad Boots        & 47  & 34  & 40  & 22 & 22 \\
\midrule
\multicolumn{6}{l}{\textit{Weapons}} \\
\quad One-handed   & 125 & 90  & 108 & 59 & 59 \\
\quad Two-handed   & 251 & 179 & 215 & 118 & 118 \\
\midrule
\multicolumn{6}{l}{\textit{Trinkets}} \\
\quad Amulet       & 157 & 108 & 133 & 71 & 72 \\
\quad Ring ($\times 2$)  & 126 & 85  & 106 & 56 & 57 \\
\quad Accessory ($\times 2$) & 110 & 74 & 92 & 49 & 50 \\
\quad Back piece   & 63  & 40  & 52  & 27 & 28 \\
\bottomrule
\end{tabular}
\end{center}

For example, a Berserker's coat (3-stat: Power major, Precision and Ferocity
minor) contributes 141~Power, 101~Precision, and 101~Ferocity.  A Marauder coat
(4-stat: Power and Precision major, Vitality and Ferocity minor) contributes
121~Power, 121~Precision, 67~Vitality, and 67~Ferocity.

\subsection{Formalising in Lean 4}

\begin{leanbox}[title={Lean 4 --- Core Data Model}]
\begin{minted}{lean}
inductive StatType where
  | power | precision | ferocity | toughness | vitality
  | conditionDamage | expertise | concentration | healingPower
  deriving DecidableEq, Repr, Hashable

inductive GearSlot where
  | headgear | shoulders | coat | gloves | leggings | boots
  | mainhand | offhand    -- or combined as twoHand
  | amulet | ring1 | ring2 | accessory1 | accessory2 | backPiece
  deriving DecidableEq, Repr

inductive ComboCategory where
  | triple   -- 1 major, 2 minor
  | quad     -- 2 major, 2 minor
  | celestial -- 7 equal
  deriving DecidableEq, Repr

structure StatCombo where
  name     : String
  category : ComboCategory
  majors   : List StatType
  minors   : List StatType
  deriving DecidableEq, Repr
\end{minted}
\end{leanbox}

The stat contribution of a combo on a given slot is determined by which
category the combo belongs to and whether a given stat is major, minor, or
(for Celestial) one of the seven included stats:

\begin{leanbox}[title={Lean 4 --- Stat Coefficients}]
\begin{minted}{lean}
structure SlotCoefficients where
  tripleMajor  : Nat
  tripleMinor  : Nat
  quadMajor    : Nat
  quadMinor    : Nat
  celestial    : Nat
  deriving Repr

def coefficients : GearSlot → SlotCoefficients
  | .headgear   => ⟨63,  45,  54,  30, 30⟩
  | .shoulders  => ⟨47,  34,  40,  22, 22⟩
  | .coat       => ⟨141, 101, 121, 67, 67⟩
  | .gloves     => ⟨47,  34,  40,  22, 22⟩
  | .leggings   => ⟨94,  67,  81,  44, 44⟩
  | .boots      => ⟨47,  34,  40,  22, 22⟩
  | .mainhand   => ⟨125, 90,  108, 59, 59⟩  -- 1H
  | .offhand    => ⟨125, 90,  108, 59, 59⟩  -- 1H
  | .amulet     => ⟨157, 108, 133, 71, 72⟩
  | .ring1      => ⟨126, 85,  106, 56, 57⟩
  | .ring2      => ⟨126, 85,  106, 56, 57⟩
  | .accessory1 => ⟨110, 74,  92,  49, 50⟩
  | .accessory2 => ⟨110, 74,  92,  49, 50⟩
  | .backPiece  => ⟨63,  40,  52,  27, 28⟩

/-- Contribution of `combo` on `slot` to `stat`. -/
def contribution (slot : GearSlot) (combo : StatCombo)
    (stat : StatType) : Nat :=
  let c := coefficients slot
  match combo.category with
  | .triple =>
    if combo.majors.contains stat then c.tripleMajor
    else if combo.minors.contains stat then c.tripleMinor
    else 0
  | .quad =>
    if combo.majors.contains stat then c.quadMajor
    else if combo.minors.contains stat then c.quadMinor
    else 0
  | .celestial =>
    let wvwStats : List StatType :=
      [.power, .precision, .ferocity, .toughness,
       .vitality, .conditionDamage, .healingPower]
    if wvwStats.contains stat then c.celestial else 0
\end{minted}
\end{leanbox}

A \emph{build} assigns a stat combo to each slot, a rune, and a distribution
of infusions:

\begin{leanbox}[title={Lean 4 --- Build and Stat Total}]
\begin{minted}{lean}
structure Build where
  gear      : GearSlot → StatCombo
  rune      : Rune   -- to be defined later
  infusions : StatType → Nat
  -- Invariant: ∑ infusions = 18

/-- Total stat from gear + infusions (excluding base stats and buffs). -/
def statTotal (b : Build) (s : StatType) : Nat :=
  let gearTotal := (GearSlot.all.map fun slot =>
    contribution slot (b.gear slot) s).sum
  gearTotal + 5 * b.infusions s
\end{minted}
\end{leanbox}

\subsection{The Objective}

Given stat targets $T : \texttt{StatType} \to \N$, we define the penalty:

\begin{defbox}[title={Asymmetric Penalty}]
\[
  \pen(b) \;=\; \sum_{s \in \texttt{StatType}}
    w_{\text{miss}} \cdot \max\!\big(0,\; T(s) - \statTotal(b, s)\big)
    \;+\;
    w_{\text{over}} \cdot \max\!\big(0,\; \statTotal(b, s) - T(s)\big)
\]
where $w_{\text{miss}} \gg w_{\text{over}}$ (e.g.\ $10 : 1$).  Missing a
target by 100 costs ten times as much as exceeding it by 100.
\end{defbox}

\begin{leanbox}[title={Lean 4 --- Penalty}]
\begin{minted}{lean}
def penalty (target actual : Nat) (wMiss wOver : Nat) : Nat :=
  if actual < target then
    wMiss * (target - actual)
  else
    wOver * (actual - target)

def totalPenalty (b : Build) (targets : StatType → Nat)
    (wMiss wOver : Nat) : Nat :=
  (StatType.all.map fun s =>
    penalty (targets s) (statTotal b s) wMiss wOver).sum
\end{minted}
\end{leanbox}

The optimisation problem: find the build $b^*$ that minimises
$\pen(b^*)$ subject to the rune constraint (all six armour runes identical),
the weapon constraint (2H or 1H+1H, not a mix), and the infusion budget
($\sum_s \texttt{infusions}(s) = 18$).

\subsection{The Search Space}

With 40 stat combos and approximately 16 gear slots, the raw search space
is $40^{16} \approx 1.2 \times 10^{25}$ gear assignments---before
considering rune choices, sigils, and infusion distributions.  Brute force is
out of the question.

\section{The Propagator Attempt}

Having read \emph{Verified Order Theory}, you recognised the structure: each
gear slot is a finite-domain variable, the rune constraint couples six of them,
and the stat targets impose global arithmetic constraints.  This is a
\emph{constraint satisfaction problem}, and you built a propagator network
to attack it.

\subsection{The Domain-Reduction Lattice}

\begin{crossrefbox}
In the lattice book, we defined a \emph{cell} as a value in a
join-semilattice, updated by joining new information.  The set-of-possible-values
lattice ordered by reverse inclusion appeared as Example~(d) in the
chapter on propagator network design.  We use that exact lattice here.
\end{crossrefbox}

For each gear slot $i$, define the \emph{domain cell} $D_i$ as a subset of
the 40 WvW-legal stat combos, ordered by reverse inclusion:
\[
  D_i \subseteq D_j \iff D_j \subseteq D_i
  \qquad\text{(fewer remaining options = more information)}.
\]
The initial value is the full set of 40 combos ($\bot$ = ``nothing ruled out'').
The join (merge) is intersection: $D_i \sqcup D_j = D_i \cap D_j$
(keep only combos consistent with both sources of information).  The top element
$\top$ is the empty set: contradiction (all combos ruled out).

The product lattice $\prod_i \Finset(\texttt{StatCombo})^{\mathrm{op}}$
represents the full network state.

\subsection{Propagators You Built}

\paragraph{The rune propagator.}
If any rune is eliminated from any armour slot's rune options, eliminate it
from all six.  If any rune is fixed (the domain becomes a singleton), fix all
six.  This is a global \texttt{AllEqual} constraint.

\paragraph{The stat-floor propagator.}
For each stat type $s$ and each unfilled slot $i$, compute
\[
  \maxContrib(i, s) = \max\{\texttt{contribution}(i, c, s) \mid c \in D_i\}.
\]
If $\sum_i \maxContrib(i, s) + 90 < T(s)$ (even with all infusions devoted
to $s$), the problem is infeasible: signal $\top$.

\paragraph{Running to fixpoint.}
Fire all propagators repeatedly until no domain changes.  By the
Knaster--Tarski theorem (proven in the lattice book), this converges to the
least fixed point of the propagator system---the maximal amount of pruning
achievable by this set of propagators.

\begin{leanbox}[title={Lean 4 --- Propagator Network (from the Lattice Book)}]
\begin{minted}{lean}
-- Reusing infrastructure from Verified Order Theory
structure Cell (L : Type) [SemilatticeSup L] where
  value : L

def Cell.update [SemilatticeSup L] (c : Cell L) (new : L) : Cell L :=
  ⟨c.value ⊔ new⟩

-- Domain cell: Finset StatCombo with reverse-inclusion order
-- ⊔ is ∩ (intersection), ⊥ is allCombos, ⊤ is ∅

-- The propagator loop from the lattice book:
partial def runToFixpoint (network : PropNetwork) : PropNetwork :=
  let network' := network.fireAll
  if network' == network then network'
  else runToFixpoint network'
\end{minted}
\end{leanbox}

\subsection{Where It Gets Stuck}

The propagator network does its job: domains shrink, impossible combos are
eliminated, and infeasible problems are detected early.  But after convergence,
you're left with \emph{pruned domains}---perhaps 8--15 combos remaining per
slot instead of 40.  That's still $15^{16} \approx 6.5 \times 10^{18}$
possible builds.

Three gaps remain:

\begin{enumerate}[nosep]
  \item \textbf{Propagation doesn't search.}  It tells you what's
    \emph{impossible}, not what's \emph{best}.  After fixpoint, multiple combos
    remain per slot, and propagation has no mechanism for choosing among them.

  \item \textbf{Propagation doesn't optimise.}  It can detect infeasibility
    (a domain goes empty, signalling $\top$) but has no notion of ``this build
    scores lower penalty than that one.''  There is no \emph{objective function}
    in the propagator framework.

  \item \textbf{Propagation doesn't handle tradeoffs.}  Should you sacrifice
    20~Power to gain 50~Vitality?  That depends on the penalty weights and
    how close you are to each target.  Propagation doesn't model comparative
    reasoning.
\end{enumerate}

\begin{intuitionbox}
Think of propagation as a \emph{filter}: it removes options that provably
can't be part of any good solution.  A \emph{solver} is a filter plus a
\emph{navigator}: it explores the remaining options systematically, tracking
the best solution found so far, and proving (by exhaustion or bounding) that
no unexplored option can beat it.
\end{intuitionbox}

The rest of this book fills these three gaps.  Chapter~\ref{ch:landscape} maps
the optimisation landscape.  Chapter~\ref{ch:stronger} makes propagation
itself stronger (so the filter removes more).  Parts~II and~III add the
navigator.  Part~IV assembles everything into a verified solver.

% ────────────────────────────────────────────────────────────────────────────
\chapter{The Landscape of Optimisation}
\label{ch:landscape}
% ────────────────────────────────────────────────────────────────────────────

\epigraph{All happy optimisation problems are alike;
each unhappy one is unhappy in its own way.}{}

Now that we have a concrete problem and understand what propagation gives us
(and doesn't), let's place the build editor in the broader landscape of
optimisation.

\section{What Is an Optimisation Problem?}

\begin{defbox}[title={Optimisation Problem}]
An \emph{optimisation problem} consists of:
\begin{itemize}[nosep]
  \item A set of \emph{decision variables} $x_1, \ldots, x_n$.
  \item For each variable $x_i$, a \emph{domain} $D_i$ of possible values.
  \item A set of \emph{constraints} $C$ that restrict which combinations
    of values are allowed.
  \item An \emph{objective function} $f : D_1 \times \cdots \times D_n \to \R$
    to minimise (or maximise).
\end{itemize}
A \emph{feasible solution} is an assignment $x_i \in D_i$ satisfying all
constraints.  An \emph{optimal solution} is a feasible solution minimising $f$.
\end{defbox}

\begin{leanbox}[title={Lean 4 --- OptProblem}]
\begin{minted}{lean}
structure OptProblem (V : Type) (D : V → Type) where
  constraints : (∀ v, D v) → Prop
  objective   : (∀ v, D v) → Int

structure Solution (P : OptProblem V D) where
  assignment : ∀ v, D v
  feasible   : P.constraints assignment

def isOptimal (P : OptProblem V D) (sol : Solution P) : Prop :=
  ∀ other : Solution P, P.objective sol.assignment ≤ P.objective other.assignment
\end{minted}
\end{leanbox}

The build editor is an instance: variables are gear slots, domains are
stat combos, constraints are the rune and weapon and infusion rules, and
the objective is the asymmetric penalty.

\section{Continuous vs.\ Discrete}

In \emph{continuous} optimisation, domains are subsets of $\R^n$.  Gradient
descent, Newton's method, and convex optimisation live here.

In \emph{discrete} (or \emph{combinatorial}) optimisation, domains are finite
sets.  The build editor is discrete: you can't choose ``60\% Berserker +
40\% Marauder''---you pick one combo per slot.  This means:

\begin{itemize}[nosep]
  \item \textbf{No gradients.}  The objective is not differentiable (it's
    defined on a finite set).  Gradient descent doesn't apply.
  \item \textbf{Local optima everywhere.}  Unlike convex problems, where every
    local minimum is global, discrete problems can have many local optima
    separated by barriers.
  \item \textbf{NP-hardness lurks.}  Many discrete optimisation problems are
    NP-hard.  We don't expect polynomial-time exact algorithms.
\end{itemize}

\section{Convexity: Why It's Not Enough}

\begin{defbox}[title={Convex Set}]
A set $S \subseteq \R^n$ is \emph{convex} if for all $x, y \in S$ and
$\lambda \in [0, 1]$: $\lambda x + (1-\lambda)y \in S$.
\end{defbox}

\begin{defbox}[title={Convex Function}]
A function $f : S \to \R$ on a convex set $S$ is \emph{convex} if for all
$x, y \in S$ and $\lambda \in [0, 1]$:
\[
  f\!\big(\lambda x + (1-\lambda)y\big) \;\leq\; \lambda f(x) + (1-\lambda) f(y).
\]
\end{defbox}

\begin{thmbox}[title={Convex $\Rightarrow$ Local = Global}]
If $f$ is convex and $x^*$ is a local minimum of $f$, then $x^*$ is a
\emph{global} minimum.
\end{thmbox}

\begin{proof}
Suppose $x^*$ is a local minimum but not global: there exists $y$ with
$f(y) < f(x^*)$.  By convexity, for small $\lambda > 0$:
\[
  f\!\big((1-\lambda)x^* + \lambda y\big)
  \leq (1-\lambda) f(x^*) + \lambda f(y)
  < f(x^*).
\]
But $(1-\lambda)x^* + \lambda y$ is arbitrarily close to $x^*$, contradicting
the local minimality of $x^*$.
\end{proof}

This is why gradient descent works for convex problems: you can't get trapped
in a local minimum.  But the build editor's domain is a finite set of points,
not a convex region of $\R^n$.  The concept of ``moving a small distance
toward a better solution'' doesn't apply.

\begin{intuitionbox}
We'll exploit convexity indirectly: by \emph{relaxing} the discrete problem
to a continuous one (pretending we can choose fractional amounts of each
stat combo), solving the relaxation, and using the solution as a \emph{bound}
on how good any discrete solution can be.  This is the LP relaxation idea
of Chapter~\ref{ch:lp}.
\end{intuitionbox}

\section{The Solver Architecture}

Here is the plan.  The gear optimiser is a stack of four layers:

\begin{center}
\begin{tikzpicture}[
  layer/.style={draw=midblue, fill=defbg, rounded corners=4pt,
    minimum width=10cm, minimum height=1.1cm, font=\small},
  arr/.style={-{Stealth[length=6pt]}, thick, midblue}
]
\node[layer] (L1) at (0, 0)
  {\textbf{Layer 1: Propagation} --- prune impossible combos (lattice book)};
\node[layer] (L2) at (0, 1.6)
  {\textbf{Layer 2: Search} --- explore remaining options systematically};
\node[layer] (L3) at (0, 3.2)
  {\textbf{Layer 3: Bounding} --- skip branches that can't beat the incumbent};
\node[layer] (L4) at (0, 4.8)
  {\textbf{Layer 4: Heuristics} --- find good starting solutions fast};
\draw[arr] (L1) -- (L2);
\draw[arr] (L2) -- (L3);
\draw[arr] (L3) -- (L4);
\end{tikzpicture}
\end{center}

Each layer wraps around the ones below it.  Propagation (Layer~1) is the
\emph{foundation}---it runs at every node of the search tree.  Search (Layer~2)
explores the pruned space.  Bounding (Layer~3) uses LP relaxation to prune
entire subtrees.  Heuristics (Layer~4) provide a warm start so bounding can
prune more aggressively.

\begin{crossrefbox}
Layer~1 is precisely the propagator network from \emph{Verified Order
Theory}.  The contribution of this book is Layers~2--4.  But notice:
propagation isn't just the foundation---it runs \emph{inside} every
search node, making the entire solver dramatically faster.
\end{crossrefbox}

% ────────────────────────────────────────────────────────────────────────────
\chapter{Making Propagation Stronger}
\label{ch:stronger}
% ────────────────────────────────────────────────────────────────────────────

\epigraph{Before adding new machinery, squeeze more out of what you have.}{}

Before we introduce search and bounding, let's strengthen the propagator
engine itself.  The propagators from the lattice book (rune equality,
basic stat-floor) are correct but leave information on the table.  This
chapter adds three new propagators that prune significantly more.

\section{The LinearGeq Propagator}

The basic stat-floor propagator detects infeasibility when the total
maximum possible contribution falls below a target.  The \texttt{LinearGeq}
propagator goes further: it prunes individual combos from individual slots.

\begin{defbox}[title={LinearGeq Propagator}]
For each stat type $s$ and each slot $i$, and for each combo $c \in D_i$:
compute the maximum possible total stat from all \emph{other} slots plus
$c$'s contribution from slot $i$ plus the full infusion budget:
\[
  U(i, c, s) = \texttt{contribution}(i, c, s) + \sum_{j \neq i} \maxContrib(j, s) + 90.
\]
If $U(i, c, s) < T(s)$, then no build using combo $c$ on slot $i$ can
meet target $T(s)$---even if every other slot takes the best possible
combo for $s$ and all infusions go to $s$.  Remove $c$ from $D_i$.
\end{defbox}

\begin{leanbox}[title={Lean 4 --- LinearGeq Propagator}]
\begin{minted}{lean}
/-- Remove combos from slot i that can't contribute enough to stat s. -/
def linearGeqPropagate
    (domains : GearSlot → Finset StatCombo)
    (targets : StatType → Nat)
    (slot : GearSlot) : Finset StatCombo :=
  domains slot |>.filter fun combo =>
    StatType.all.all fun stat =>
      let myContrib := contribution slot combo stat
      let othersMax := (GearSlot.all.filter (· ≠ slot) |>.map fun j =>
        maxContribForStat (domains j) j stat).sum
      let infusionBudget := 90  -- 18 slots × 5 each
      myContrib + othersMax + infusionBudget ≥ targets stat
\end{minted}
\end{leanbox}

\begin{propbox}[title={Monotonicity of LinearGeq}]
The \texttt{LinearGeq} propagator is monotone: if domains shrink
(more information), the propagator can only remove \emph{more} combos
(produce more information), never fewer.
\end{propbox}

\begin{proof}
Suppose $D'_j \subseteq D_j$ for all $j$ (domains have shrunk in the
reverse-inclusion order, meaning we have more information).  For any slot $i$
and combo $c$:
\[
  \sum_{j \neq i} \maxContrib'(j, s)
  \;\leq\;
  \sum_{j \neq i} \maxContrib(j, s)
\]
because $\maxContrib$ can only decrease when the domain shrinks.  So $U'(i,c,s)
\leq U(i,c,s)$.  Any combo removed by the old propagator is still removed by the
new one (with smaller $U$), and possibly additional combos are removed.  Hence
the output domain $D'_i$ under the new inputs is a subset of the output
under the old inputs---i.e., the propagator is monotone in the reverse-inclusion
order.
\end{proof}

\section{The LinearLeq Propagator}

The dual of \texttt{LinearGeq}: if we want to minimise overshoot, we can
prune combos that would cause unavoidable excess.

\begin{defbox}[title={LinearLeq Propagator}]
For each stat $s$, slot $i$, and combo $c \in D_i$: compute
\[
  L(i, c, s) = \texttt{contribution}(i, c, s) + \sum_{j \neq i} \minContrib(j, s).
\]
If $L(i, c, s) > T(s) + \text{maxAllowedOvershoot}$, remove $c$ from $D_i$.
This is weaker (since the overshoot penalty weight is low), but still helps
when a stat is close to its target.
\end{defbox}

\begin{leanbox}[title={Lean 4 --- LinearLeq Propagator}]
\begin{minted}{lean}
/-- Remove combos from slot i that would force unavoidable overshoot
    on some stat beyond the allowed threshold. -/
def linearLeqPropagate
    (domains : GearSlot → Finset StatCombo)
    (targets : StatType → Nat)
    (maxOvershoot : Nat)  -- e.g., 200 stat points
    (slot : GearSlot) : Finset StatCombo :=
  domains slot |>.filter fun combo =>
    StatType.all.all fun stat =>
      let myContrib := contribution slot combo stat
      let othersMin := (GearSlot.all.filter (· ≠ slot) |>.map fun j =>
        minContribForStat (domains j) j stat).sum
      myContrib + othersMin ≤ targets stat + maxOvershoot
\end{minted}
\end{leanbox}

\begin{propbox}[title={Monotonicity of LinearLeq}]
\texttt{LinearLeq} is monotone: if domains shrink, the propagator removes
at least as many combos.
\end{propbox}

\begin{proof}
When domains shrink ($D'_j \subseteq D_j$), the minimum contribution from
other slots can only \emph{increase} (the minimum over a smaller set is
$\geq$ the minimum over a larger set):
$\minContrib'(j, s) \geq \minContrib(j, s)$.
So $L'(i, c, s) \geq L(i, c, s)$.  Any combo removed by the old propagator
(because $L(i, c, s) > T(s) + \text{threshold}$) is still removed.
\end{proof}

Note that \texttt{LinearLeq} is most useful for stats where the player
explicitly \emph{doesn't want} too much.  In WvW, this might occur when
a player wants to maximise damage stats while keeping Toughness below a
threshold (to avoid drawing aggro from certain AI enemies that target the
highest-Toughness player).

\section{The Infusion Budget Propagator}

\begin{defbox}[title={Infusion Budget Propagator}]
After all gear propagation has reached fixpoint, compute, for each stat $s$:
\[
  \text{bestGear}(s) = \sum_i \maxContrib(i, s).
\]
The \emph{deficit} for $s$ is $\max(0, T(s) - \text{bestGear}(s))$.
If $\sum_s \text{deficit}(s) > 90$ (the total deficit across all stats
exceeds the infusion budget), the problem is infeasible.
\end{defbox}

This is a global constraint that reasons about the interaction between gear
choices and the infusion budget.  It can detect infeasibility that no
single-stat propagator would catch: each stat's deficit might be small
enough to fill, but the total across all stats exceeds the budget.

\begin{leanbox}[title={Lean 4 --- Infusion Budget Propagator}]
\begin{minted}{lean}
/-- Check whether the total stat deficit across all stats exceeds
    the infusion budget (18 slots × 5 = 90 points). -/
def infusionBudgetFeasible
    (domains : GearSlot → Finset StatCombo)
    (targets : StatType → Nat) : Bool :=
  let totalDeficit := StatType.all.map (fun stat =>
    let bestGear := (GearSlot.all.map fun slot =>
      maxContribForStat (domains slot) slot stat).sum
    if targets stat > bestGear then targets stat - bestGear
    else 0) |>.sum
  totalDeficit ≤ 90
\end{minted}
\end{leanbox}

\section{The Strengthened Network}

Combining all propagators into a single network:

\begin{enumerate}[nosep]
  \item Fire the rune propagator (from the lattice book).
  \item Fire \texttt{LinearGeq} for every slot and stat.
  \item Fire \texttt{LinearLeq} for every slot and stat (if overshoot control
    is enabled).
  \item Check the infusion budget propagator.
  \item Repeat until fixpoint.
\end{enumerate}

\begin{thmbox}[title={Termination of the Strengthened Network}]
The strengthened propagator network terminates in at most
$\sum_i |D_i^{(0)}|$ steps, where $D_i^{(0)}$ is the initial domain
of slot $i$.  With 40 combos and 16 slots, this is at most 640 steps.
\end{thmbox}

\begin{proof}
Each propagation step either removes at least one combo from at least one
domain (strictly decreasing the lattice height) or is a no-op (fixpoint
reached).  The lattice height is $\sum_i |D_i|$, initially at most
$40 \times 16 = 640$.  Since the height is a natural number that strictly
decreases at each non-trivial step, the process terminates in at most
640 steps.
\end{proof}

\begin{exercisebox}
Implement the full strengthened propagator network in Lean~4, reusing the
\texttt{Cell} and \texttt{Propagator} types from the lattice book.
Run it on the following problem: Heavy armour class, stat targets
Power $\geq 2500$, Precision $\geq 2000$, Ferocity $\geq 1200$.
How many combos remain per slot after propagation?
\end{exercisebox}

\begin{exercisebox}
Prove in Lean~4 that the \texttt{LinearGeq} propagator satisfies the
monotonicity property from the lattice book: if input domains shrink,
output domains shrink.
\end{exercisebox}

% ════════════════════════════════════════════════════════════════════════════
%  PART II — THE MATHEMATICS OF OPTIMISATION
% ════════════════════════════════════════════════════════════════════════════
\part{The Mathematics of Optimisation}

% ────────────────────────────────────────────────────────────────────────────
\chapter{Linear Programming}
\label{ch:lp}
% ────────────────────────────────────────────────────────────────────────────

\epigraph{The shortest distance between two points is a straight line,\\
and the best solution to a linear problem is at a corner.}{}

\section{The LP Standard Form}

\begin{defbox}[title={Linear Program (Standard Form)}]
\[
  \min_{x \in \R^n} \; c^\top x
  \quad\text{subject to}\quad
  Ax \leq b, \quad x \geq 0.
\]
Here $c \in \R^n$ is the cost vector, $A \in \R^{m \times n}$ is the
constraint matrix, and $b \in \R^m$ is the right-hand side.
\end{defbox}

The feasible region $\{x \mid Ax \leq b, x \geq 0\}$ is a convex
\emph{polytope}---the intersection of finitely many half-spaces.

\begin{thmbox}[title={Fundamental Theorem of Linear Programming}]
If an LP has an optimal solution, then it has an optimal solution at a
\emph{vertex} (extreme point) of the feasible polytope.
\end{thmbox}

\begin{intuitionbox}
Why?  The objective $c^\top x$ is a linear function.  On any edge of the
polytope, it is monotonically increasing or decreasing.  So moving along
an edge toward a vertex can only improve (or maintain) the objective.
The optimum must occur where you can't improve further---at a vertex.
\end{intuitionbox}

\section{The Simplex Method}

The simplex method exploits this: start at a vertex, move to an adjacent
vertex with a better objective value, repeat until no improving neighbour
exists.

\begin{crossrefbox}
The vertices of the feasible polytope, ordered by objective value,
form a partial order.  The simplex method is a \emph{monotone walk}
on this structure---the same pattern as propagator convergence.  In the
lattice book, cell values only moved upward; here, the objective only
moves downward (since we're minimising).  The fixed point is the optimum.
\end{crossrefbox}

\subsection{Pivoting}

A vertex of the polytope corresponds to a \emph{basic feasible solution}
(BFS): a choice of $m$ constraints to hold with equality (the \emph{basis}).
The simplex method maintains a basis and performs \emph{pivot} operations
to swap one constraint in and one out, moving to an adjacent vertex.

The pivot rule selects which non-basic variable to bring into the basis
(improving the objective) and which basic variable to remove (maintaining
feasibility).  Bland's rule (choose the smallest-index improving variable)
guarantees termination by preventing cycling.

\begin{leanbox}[title={Lean 4 --- Simplex Sketch}]
\begin{minted}{lean}
structure Tableau where
  m : Nat  -- number of constraints
  n : Nat  -- number of variables
  tab : Array (Array Float)  -- (m+1) × (n+m+1) augmented matrix
  basis : Array Nat          -- indices of basic variables

/-- One pivot step: bring variable enterIdx into the basis. -/
def pivot (T : Tableau) (enterIdx : Nat) : Option Tableau :=
  -- Minimum ratio test to find the leaving variable
  let leaveIdx? := minRatioTest T enterIdx
  match leaveIdx? with
  | none => none  -- unbounded
  | some leaveIdx =>
    -- Gaussian elimination to pivot
    some (gaussianPivot T enterIdx leaveIdx)

/-- Run simplex to completion. -/
partial def simplex (T : Tableau) : SimplexResult :=
  match blandEnteringVariable T with
  | none => .optimal (extractSolution T)  -- no improving variable
  | some enterIdx =>
    match pivot T enterIdx with
    | none => .unbounded
    | some T' => simplex T'
\end{minted}
\end{leanbox}

\section{Duality}

Every LP has a \emph{dual} problem.  If the primal is
$\min\; c^\top x$ s.t.\ $Ax \leq b, x \geq 0$,
the dual is $\max\; b^\top y$ s.t.\ $A^\top y \geq c, y \geq 0$.

\begin{thmbox}[title={Weak Duality}]
For any primal feasible $x$ and dual feasible $y$:
$c^\top x \geq b^\top y$.
\end{thmbox}

\begin{proof}
$c^\top x \leq (A^\top y)^\top x = y^\top A x \leq y^\top b = b^\top y$.
The first inequality uses dual feasibility ($A^\top y \geq c$ and $x \geq 0$);
the second uses primal feasibility ($Ax \leq b$ and $y \geq 0$).
\end{proof}

\begin{thmbox}[title={Strong Duality}]
If the primal has an optimal solution $x^*$, then the dual has an optimal
solution $y^*$ with $c^\top x^* = b^\top y^*$.
\end{thmbox}

The dual provides \emph{certificates of optimality}: if you exhibit primal and
dual solutions with matching objective values, you've \emph{proved} the primal
solution is optimal, without trusting the simplex algorithm's internal state.

\begin{intuitionbox}
In the gear solver, the LP relaxation gives a \emph{lower bound} on the
achievable penalty.  If we also produce the dual solution, we have a
\emph{machine-checkable certificate} that no gear loadout can do better
than this bound.  When the integer solution matches the LP bound,
optimality is proven.
\end{intuitionbox}

\section{LP Relaxation of the Gear Problem}

Here is how to relax the build editor to an LP.  For each slot $i$ and
combo $c$, introduce a variable $x_{i,c} \in [0, 1]$ representing the
``fraction'' of combo $c$ used on slot $i$.  The constraints:

\begin{enumerate}[nosep]
  \item $\sum_c x_{i,c} = 1$ for each slot $i$ (exactly one combo per slot).
  \item $x_{i,c} \geq 0$ for all $i, c$.
  \item The stat contributions are linear in $x$:
    $\statTotal(x, s) = \sum_{i,c} x_{i,c} \cdot \texttt{contribution}(i, c, s)$.
  \item Linearise the penalty: introduce auxiliary variables for the
    max-with-zero terms.
\end{enumerate}

The LP relaxation drops the constraint $x_{i,c} \in \{0, 1\}$ and allows
fractional values.  The optimal LP value is a lower bound on the optimal
integer value.

\begin{exercisebox}
Formulate the LP relaxation for a small instance (3 armour slots, 5 stat
combos, 2 stat targets).  Solve it by hand or with the simplex implementation.
Verify that the LP optimum is $\leq$ the best integer solution.
\end{exercisebox}

% ────────────────────────────────────────────────────────────────────────────
\chapter{Integer Programming and Branch-and-Bound}
\label{ch:bb}
% ────────────────────────────────────────────────────────────────────────────

\epigraph{Divide and conquer, but know when to surrender.}{}

\section{The Integrality Gap}

The LP relaxation allows fractional solutions like ``use 60\% Berserker and
40\% Marauder on the coat.''  The \emph{integrality gap} is the difference
between the LP optimum and the best integer (actual gear) solution.

\begin{defbox}[title={Integrality Gap}]
\[
  \text{gap} = \opt(\IP) - \opt(\LP)
\]
where $\opt(\IP)$ is the optimal integer (actual gear) penalty and $\opt(\LP)$
is the optimal LP relaxation penalty.  Since $\opt(\LP) \leq \opt(\IP)$
(relaxation can only improve the objective), the gap is always $\geq 0$.
\end{defbox}

\begin{intuitionbox}
Consider a simplified two-slot, two-stat problem.  Slot~A can be Berserker's
(+141 Power, +0 Toughness) or Knight's (+0 Power, +141 Toughness).  Slot~B
is the same.  Targets: Power $\geq 100$, Toughness $\geq 100$.

\emph{LP relaxation:} use 50\% Berserker + 50\% Knight on each slot.  Each
slot contributes $70.5$ Power and $70.5$ Toughness.  Total: $141$ Power,
$141$ Toughness.  Both targets met.  Penalty = 0.

\emph{Best integer solution:} one slot Berserker's, one slot Knight's.
Total: $141$ Power, $141$ Toughness.  Penalty = 0.

Here the gap is zero---the LP and IP optima coincide.  But now add a third
stat target (Ferocity $\geq 50$) where Berserker's gives 101 Ferocity and
Knight's gives 0.  The LP can get Ferocity from a fractional Berserker
assignment on both slots; the IP must concentrate it on one, potentially
overshooting other targets.  The gap opens up.
\end{intuitionbox}

The gap can be large: the LP might achieve penalty~50 by artfully spreading
stat contributions across fractional combos, while the best integer build has
penalty~200 because it must commit each slot to a single combo.

\begin{histbox}
The study of integrality gaps is central to the theory of approximation
algorithms.  Nemhauser and Wolsey's \emph{Integer and Combinatorial
Optimization} (1988) established much of the foundational theory.  A small
gap means the LP relaxation gives a tight bound and B\&B converges quickly;
a large gap means the LP bound is loose and more search is needed.  For the
gear problem, the gap is typically small when targets are tight (close to
the maximum achievable) and larger when targets are loose.
\end{histbox}

\section{Branch and Bound}

Branch and bound (B\&B) is a systematic search that exploits the LP bound
to prune.

\begin{defbox}[title={Branch and Bound}]
\begin{enumerate}[nosep]
  \item \textbf{Initialise.}  Set \texttt{incumbent} (best solution found)
    to $\infty$.  Create a root node representing the full problem.
  \item \textbf{Select.}  Pick an unexplored node from the queue.
  \item \textbf{Bound.}  Solve the LP relaxation of this node's subproblem.
    If the LP optimum $\geq$ \texttt{incumbent}, \emph{prune} (no solution
    in this subtree can beat the incumbent).
  \item \textbf{Branch.}  If not pruned and the LP solution is fractional,
    choose a slot $i$ with a fractional assignment.  Create child nodes,
    one for each remaining combo in $D_i$ (fixing slot $i$ to that combo).
  \item \textbf{Propagate.}  At each child node, run the propagator network.
    If any domain becomes empty, prune.
  \item \textbf{Repeat.}  Go to step 2 until the queue is empty.
    The incumbent is optimal.
\end{enumerate}
\end{defbox}

\begin{crossrefbox}
The B\&B search tree is a \emph{lattice of subproblems}, ordered by
refinement (fixing more variables).  At each node, the LP bound is an
upper bound on the subtree's optimal value.  Pruning = detecting that
a subtree is ``dead''---analogous to detecting $\top$ (contradiction)
in a propagator cell.  The whole B\&B is a monotone process: as we
descend the tree, domains only shrink, bounds only tighten.
\end{crossrefbox}

\subsection{The B\&B Invariant}

The key correctness invariant, which we'll formalise and prove in the capstone:

\begin{thmbox}[title={B\&B Invariant}]
At every point during the B\&B search:
\begin{enumerate}[nosep]
  \item The incumbent is a \emph{feasible} solution (satisfies all constraints).
  \item Every pruned node $n$ satisfies $\LP(n) \geq \texttt{incumbent.penalty}$
    (its LP bound is no better than the incumbent).
  \item Every integer solution in a pruned subtree has penalty $\geq \LP(n)
    \geq \texttt{incumbent.penalty}$ (by weak duality and the relaxation
    relationship).
\end{enumerate}
When the queue is empty, these invariants imply the incumbent is optimal:
every feasible solution is either the incumbent itself, or lives in a
pruned subtree with penalty $\geq$ incumbent.
\end{thmbox}

\subsection{Variable and Value Ordering}

The order in which we branch on slots and try combos dramatically affects
the number of nodes explored.

\paragraph{Fail-first variable ordering.}
At each branch node, choose the slot with the \emph{smallest remaining
domain} (after propagation).  Intuition: if a slot has only 2 remaining
combos, branching on it creates 2 children.  If another has 15, branching
creates 15 children.  Fail-first focuses effort where the problem is most
constrained.

\begin{leanbox}[title={Lean 4 --- Fail-First Variable Selection}]
\begin{minted}{lean}
/-- Select the unfixed slot with the smallest remaining domain. -/
def selectSlotFailFirst
    (domains : GearSlot → Finset StatCombo) : GearSlot :=
  let unfixed := GearSlot.all.filter fun s => (domains s).card > 1
  unfixed.foldl (init := unfixed.head!) fun best slot =>
    if (domains slot).card < (domains best).card then slot else best
\end{minted}
\end{leanbox}

\paragraph{Best-first value ordering.}
Within a slot, try combos in order of their contribution to the
\emph{tightest} stat target first.  Define ``tightest'' as the stat with
the largest remaining deficit (target minus current best achievable):

\begin{leanbox}[title={Lean 4 --- Value Ordering by Tightest Deficit}]
\begin{minted}{lean}
/-- Find the stat with the largest remaining deficit. -/
def tightestStat (domains : GearSlot → Finset StatCombo)
    (targets : StatType → Nat) : StatType :=
  StatType.all.foldl (init := .power) fun best stat =>
    let bestGear := (GearSlot.all.map fun s =>
      maxContribForStat (domains s) s stat).sum
    let deficit := if targets stat > bestGear then targets stat - bestGear else 0
    let bestDeficit := if targets best > (GearSlot.all.map fun s =>
      maxContribForStat (domains s) s best).sum then targets best - (GearSlot.all.map fun s =>
      maxContribForStat (domains s) s best).sum else 0
    if deficit > bestDeficit then stat else best

/-- Order combos by contribution to the tightest stat (descending). -/
def orderCombos (slot : GearSlot) (combos : Finset StatCombo)
    (tightest : StatType) : List StatCombo :=
  combos.toList.mergeSort fun a b =>
    contribution slot a tightest > contribution slot b tightest
\end{minted}
\end{leanbox}

This increases the chance of finding a good solution early in the search,
improving the incumbent and enabling more LP-based pruning.

\subsection{Symmetry Breaking}

\begin{propbox}[title={Ring and Accessory Symmetry}]
Ring~1 and Ring~2 are interchangeable: swapping their stat combos produces
an identical build.  The same holds for Accessory~1 and Accessory~2.
\end{propbox}

Impose a total order on stat combos (e.g.\ alphabetical by name) and
require $\texttt{combo}(\texttt{ring}_1) \leq \texttt{combo}(\texttt{ring}_2)$.
This halves the search space for each symmetric pair without eliminating any
essentially different builds.

\begin{leanbox}[title={Lean 4 --- Symmetry Breaking}]
\begin{minted}{lean}
/-- A build is canonical if symmetric slot pairs are ordered. -/
def Build.isCanonical (b : Build) : Prop :=
  comboOrd (b.gear .ring1) ≤ comboOrd (b.gear .ring2) ∧
  comboOrd (b.gear .accessory1) ≤ comboOrd (b.gear .accessory2)

/-- Every build has a canonical equivalent with the same penalty. -/
theorem canonical_exists (b : Build) :
    ∃ b', b'.isCanonical ∧ totalPenalty b' = totalPenalty b := by
  -- Swap rings/accessories if needed; stat totals are unchanged
  -- since contribution depends only on the slot type, not the index
  sorry -- exercise for the reader
\end{minted}
\end{leanbox}

\begin{exercisebox}
Prove the \texttt{canonical\_exists} theorem.  Hint: you need the fact that
\texttt{contribution .ring1 = contribution .ring2} (the coefficient table
gives the same values for both ring slots).
\end{exercisebox}

% ────────────────────────────────────────────────────────────────────────────
\chapter{Search Strategies}
\label{ch:search}
% ────────────────────────────────────────────────────────────────────────────

\epigraph{Search is the art of not looking where you don't need to.}{}

\section{Backtracking Search with Propagation}

The simplest complete search strategy: choose a variable (slot), try each
value (combo), propagate, recurse.

\begin{leanbox}[title={Lean 4 --- Backtracking Search}]
\begin{minted}{lean}
/-- Backtracking search with propagation at each node. -/
partial def backtrack
    (domains : GearSlot → Finset StatCombo)
    (assigned : List (GearSlot × StatCombo))
    (best : Option (Build × Nat))  -- incumbent
    (targets : StatType → Nat) : Option (Build × Nat) :=
  -- 1. Propagate
  let domains' := propagateToFixpoint domains
  -- 2. Check for empty domain (contradiction)
  if GearSlot.all.any (fun s => (domains' s).isEmpty) then best
  -- 3. Check if all assigned
  else if GearSlot.all.all (fun s => (domains' s).card = 1) then
    let build := extractBuild domains'
    let pen := totalPenalty build targets wMiss wOver
    match best with
    | none => some (build, pen)
    | some (_, bestPen) =>
      if pen < bestPen then some (build, pen) else best
  -- 4. Branch on smallest domain (fail-first)
  else
    let slot := selectSlot domains'  -- smallest non-singleton domain
    (domains' slot).fold best fun acc combo =>
      let domains'' := fixSlot domains' slot combo
      backtrack domains'' ((slot, combo) :: assigned) acc targets
\end{minted}
\end{leanbox}

\subsection{Worked Example: A Tiny Gear Problem}

To build intuition, consider a stripped-down problem with three armour slots
(coat, leggings, boots), three stat combos (Berserker's, Marauder, Celestial),
and a single stat target: Power $\geq 300$.

The Power contributions at Ascended:
\begin{center}
\small
\begin{tabular}{lccc}
\toprule
& \textbf{Berserker's} & \textbf{Marauder} & \textbf{Celestial} \\
& (3-stat Maj) & (4-stat Maj) & (Celestial) \\
\midrule
Coat     & 141 & 121 & 67 \\
Leggings & 94  & 81  & 44 \\
Boots    & 47  & 40  & 22 \\
\bottomrule
\end{tabular}
\end{center}

\paragraph{Step 1: Propagation.}
The \texttt{LinearGeq} propagator checks each slot.  For boots with Celestial:
the maximum Power from the other two slots is $141 + 94 = 235$, plus
Celestial boots gives $22$, total $257 < 300$.  Even with infusions (say,
no infusions in this simplified version), Celestial on boots is infeasible
if we need Power $\geq 300$.  Remove it.

Similarly, Celestial on leggings gives at most $141 + 44 + 47 = 232 < 300$.
Remove.  Celestial on coat: $67 + 94 + 47 = 208 < 300$.  Remove.

After propagation: every slot has domain $\{\text{Berserker's}, \text{Marauder}\}$.
The search space dropped from $3^3 = 27$ to $2^3 = 8$.

\paragraph{Step 2: Search.}
All domains have size~2 (tie for fail-first).  Pick coat.  Try Berserker's:
$141 + \max(94, 81) + \max(47, 40) = 141 + 94 + 47 = 282 < 300$?  No---we
haven't checked the actual assignment, just the maximum.  Fix coat to
Berserker's, propagate again.  Now leggings with Marauder gives
$141 + 81 + 47 = 269 < 300$ (still checking maximums of remaining).
We need to actually evaluate: Berserker's on all three gives
$141 + 94 + 47 = 282$; Berserker/Berserker/Marauder gives
$141 + 94 + 40 = 275$; and so on.  The best is all-Berserker's at $282$,
which still falls short of $300$.

In the full problem, infusions would close this gap ($18 \times 5 = 90$
free points).  But the example illustrates the propagate-then-search pattern.

\paragraph{Step 3: Result.}
If $282 + 90 = 372 \geq 300$ (with infusion budget), the build is feasible.
The solver would continue searching to find the build with minimum total
penalty across all stats.

\section{A* and Best-First Search}

Instead of depth-first search (which explores deeply before backtracking),
best-first search maintains a priority queue of nodes ordered by their LP
bound.  It always expands the most promising node first.

\begin{defbox}[title={A* for Optimisation}]
Maintain a priority queue of search nodes, ordered by
$f(n) = g(n) + h(n)$ where:
\begin{itemize}[nosep]
  \item $g(n)$ = cost of the partial assignment so far (penalty from already-fixed slots).
  \item $h(n)$ = heuristic lower bound on the remaining cost (LP relaxation of unfixed slots).
\end{itemize}
Always expand the node with smallest $f(n)$.
\end{defbox}

\begin{thmbox}[title={Admissibility of LP Heuristic}]
If $h(n) \leq h^*(n)$ (the LP bound is a lower bound on the true remaining
cost), then A* finds an optimal solution.
\end{thmbox}

\begin{proof}
Suppose A* terminates by expanding a goal node $n^*$ (a complete assignment).
At that moment, every unexplored node $m$ in the queue satisfies
$f(m) \geq f(n^*)$ (A* always expands the smallest-$f$ node).
For any unexplored node $m$: $f(m) = g(m) + h(m) \leq g(m) + h^*(m)$
by admissibility.  But $g(m) + h^*(m)$ is the best objective achievable
through $m$.  So the best objective through any unexplored path is $\geq
f(n^*) = g(n^*)$ (since $h(n^*) = 0$ for a complete assignment).
Therefore $n^*$ is optimal.
\end{proof}

\begin{leanbox}[title={Lean 4 --- A* Search Skeleton}]
\begin{minted}{lean}
structure SearchNode where
  domains  : GearSlot → Finset StatCombo
  assigned : List (GearSlot × StatCombo)
  gCost    : Nat   -- penalty from fixed slots
  hCost    : Nat   -- LP bound on remaining
  deriving Repr

def SearchNode.fCost (n : SearchNode) : Nat := n.gCost + n.hCost

/-- A* search for the optimal gear build. -/
partial def aStarSearch
    (queue : BinaryHeap SearchNode (·.fCost < ·.fCost))
    (incumbent : Option (Build × Nat))
    (targets : StatType → Nat) : Option (Build × Nat) :=
  match queue.extractMin with
  | none => incumbent  -- queue empty: incumbent is optimal
  | some (node, queue') =>
    -- Pruning: if f(node) ≥ incumbent, skip
    match incumbent with
    | some (_, bestPen) =>
      if node.fCost ≥ bestPen then aStarSearch queue' incumbent targets
      else expandNode node queue' incumbent targets
    | none => expandNode node queue' incumbent targets
\end{minted}
\end{leanbox}

This is a standard result from AI search theory.  The LP relaxation satisfies
admissibility because relaxing integrality constraints can only improve
(lower) the objective.

\begin{histbox}
A* was invented by Peter Hart, Nils Nilsson, and Bertram Raphael in 1968
at the Stanford Research Institute.  The ``A'' stands for ``Algorithm''
(it was called Algorithm~A with admissible heuristics, hence A*).  It
remains the gold standard for optimal search with heuristics.  In our
solver, A* with LP-relaxation heuristics is the combination of Layers~2
and~3 (search + bounding) from the architecture diagram.
\end{histbox}

\section{Nogood Learning}

When backtracking from a dead end, we can record the combination of
assignments that caused the failure and prevent future exploration of any
branch that contains the same combination.

\begin{defbox}[title={Nogood}]
A \emph{nogood} is a partial assignment that is provably inconsistent.
If $\{\text{coat} \mapsto \text{Celestial}, \text{leggings} \mapsto
\text{Celestial}\}$ causes propagation to detect infeasibility (Power
target unachievable), then this pair is a nogood.  Any future branch
containing both assignments can be pruned immediately.
\end{defbox}

\begin{leanbox}[title={Lean 4 --- Nogood Store}]
\begin{minted}{lean}
/-- A nogood is a set of (slot, combo) assignments known to be infeasible. -/
structure Nogood where
  assignments : List (GearSlot × StatCombo)

/-- Check if the current assignment contains any known nogood. -/
def isNogoodHit (nogoods : List Nogood) (assigned : List (GearSlot × StatCombo))
    : Bool :=
  nogoods.any fun ng =>
    ng.assignments.all fun (slot, combo) =>
      assigned.any fun (s, c) => s == slot && c == combo

/-- Extract a nogood from a failed propagation.
    Analyse which assignments contributed to the empty domain. -/
def extractNogood (assigned : List (GearSlot × StatCombo))
    (failedSlot : GearSlot) (failedStat : StatType)
    : Nogood :=
  -- Keep only assignments whose contribution to failedStat
  -- was necessary for the failure
  let relevant := assigned.filter fun (slot, combo) =>
    contribution slot combo failedStat > 0
  ⟨relevant⟩
\end{minted}
\end{leanbox}

\begin{crossrefbox}
This is the same pattern as CDCL (conflict-driven clause learning)
in SAT solving, which we discussed in the lattice book as ``adding
a new propagator at runtime.''  A nogood is a constraint learned from
a conflict; adding it to the network prevents the same conflict from
recurring.  In the lattice-theoretic view, a nogood is a new propagator
that removes assignments from the product lattice of domain cells.
\end{crossrefbox}

\begin{exercisebox}
Extend the backtracking search from earlier in this chapter to maintain
a nogood store.  When propagation fails, extract a nogood and add it.
Before exploring each branch, check the nogood store.  Measure the
reduction in explored nodes on a test instance.
\end{exercisebox}

\begin{exercisebox}
Prove in Lean~4 that nogood learning preserves completeness: if a build
was optimal before learning, it is still reachable after learning.  Hint:
nogoods only prune infeasible partial assignments.
\end{exercisebox}

% ────────────────────────────────────────────────────────────────────────────
\chapter{Relaxation and Bounding}
\label{ch:relax}
% ────────────────────────────────────────────────────────────────────────────

\epigraph{To understand the integer world,\\first solve the fractional one.}{}

\section{Lagrangian Relaxation}

Instead of relaxing integrality, we can relax \emph{constraints}.  The rune
constraint (all six armour rune slots must match) is a coupling constraint
that prevents us from optimising each slot independently.

\begin{defbox}[title={Lagrangian Relaxation}]
Move the rune constraint into the objective as a penalty:
\[
  L(\lambda) = \min_x \; \pen(x) + \sum_{i < j}
    \lambda_{ij} \cdot \mathbb{1}[\text{rune}(i) \neq \text{rune}(j)]
\]
where $\lambda_{ij} \geq 0$ are \emph{Lagrange multipliers}.  For any
$\lambda \geq 0$, $L(\lambda)$ is a lower bound on the optimal penalty.
\end{defbox}

\begin{thmbox}[title={Weak Lagrangian Duality}]
For any $\lambda \geq 0$: $L(\lambda) \leq \opt(\text{original problem})$.
\end{thmbox}

\begin{proof}
Let $x^*$ be an optimal solution to the original problem.  Since $x^*$
satisfies the rune constraint, $\mathbb{1}[\text{rune}(i) \neq
\text{rune}(j)] = 0$ for all $i, j$ in $x^*$.  Therefore:
\[
  L(\lambda) \leq \pen(x^*) + \sum_{i < j} \lambda_{ij} \cdot 0
  = \pen(x^*) = \opt.
\]
The inequality holds because $L(\lambda)$ is a minimum over \emph{all}
assignments (including ones violating the rune constraint), while $x^*$
is just one specific feasible assignment.
\end{proof}

With the rune constraint relaxed, the problem decomposes into independent
per-slot subproblems, each trivially solvable by enumeration (pick the
best combo for each slot independently).

\subsection{Subgradient Optimisation}

The \emph{Lagrangian dual} is $\max_{\lambda \geq 0} L(\lambda)$: find the
multipliers that give the tightest bound.  This is a non-smooth optimisation
problem (the function $L$ is piecewise linear), solved by the subgradient
method:

\begin{leanbox}[title={Lean 4 --- Subgradient Method for Lagrangian Dual}]
\begin{minted}{lean}
/-- One iteration of subgradient optimisation. -/
def subgradientStep (lambda : Array Float) (stepSize : Float)
    (solution : GearSlot → StatCombo) : Array Float :=
  -- Subgradient: violation of the rune constraint
  -- g_ij = 1 if rune(i) ≠ rune(j), else 0
  let armourSlots := [.headgear, .shoulders, .coat, .gloves, .leggings, .boots]
  let violations := armourSlots.pairs.map fun (i, j) =>
    if runeOf (solution i) == runeOf (solution j) then 0.0 else 1.0
  -- Update: λ' = max(0, λ + stepSize * g)
  lambda.zipWith violations.toArray fun li gi =>
    Float.max 0.0 (li + stepSize * gi)

/-- Run the subgradient method for a fixed number of iterations. -/
def lagrangianDual (problem : GearProblem) (iters : Nat := 100)
    : Float × (GearSlot → StatCombo) :=
  let init_lambda := Array.mkArray 15 0.0  -- C(6,2) = 15 pairs
  (List.range iters).foldl (init := (0.0, init_lambda, none)) fun (bestBound, lambda, bestSol) _ =>
    -- Solve Lagrangian subproblem (per-slot enumeration)
    let solution := solveRelaxedPerSlot lambda problem.targets
    let bound := lagrangianObjective solution lambda problem.targets
    let stepSize := 1.0 / (1.0 + iters.toFloat)  -- diminishing step
    let lambda' := subgradientStep lambda stepSize solution
    let bestBound' := Float.max bestBound bound
    (bestBound', lambda', some solution)
  |>.map fun (bound, _, sol) => (bound, sol.getD defaultSolution)
\end{minted}
\end{leanbox}

\begin{histbox}
Lagrangian relaxation was pioneered by Marshall Fisher and Michael Held
in the 1970s for vehicle routing and scheduling problems.  The key insight
is that \emph{complicating constraints}---the ones that prevent
decomposition---can be moved into the objective, trading hard constraints
for soft penalties.  The resulting decomposition often reveals beautiful
structure: in our case, the rune constraint is the only thing preventing
each slot from being optimised independently.
\end{histbox}

\section{Reduced-Cost Fixing}

\begin{defbox}[title={Reduced-Cost Fixing}]
After solving the LP relaxation, each non-basic variable $x_{i,c}$ has a
\emph{reduced cost} $\bar{c}_{i,c}$: the amount the objective would increase
if we forced $x_{i,c}$ to become positive.  If $\LP + \bar{c}_{i,c} \geq
\texttt{incumbent}$, then any solution using combo $c$ on slot $i$ has penalty
$\geq$ incumbent.  We can permanently remove $c$ from $D_i$.
\end{defbox}

This is \emph{propagation informed by relaxation}: the LP relaxation feeds
information back into the propagator network, tightening domains beyond what
pure constraint reasoning can achieve.

\begin{leanbox}[title={Lean 4 --- Reduced-Cost Fixing}]
\begin{minted}{lean}
/-- Remove combos whose reduced cost proves they can't be in
    any solution better than the incumbent. -/
def reducedCostFix
    (domains : GearSlot → Finset StatCombo)
    (lpResult : LPResult)
    (incumbentPen : Nat) : GearSlot → Finset StatCombo :=
  fun slot => (domains slot).filter fun combo =>
    let rc := lpResult.reducedCost slot combo
    -- Keep combo only if LP bound + reduced cost < incumbent
    lpResult.objectiveValue + rc < incumbentPen.toFloat
\end{minted}
\end{leanbox}

\begin{crossrefbox}
Reduced-cost fixing bridges the LP world and the propagator world.  The LP
relaxation is ``above'' the propagator network in the solver stack, but its
conclusions (``combo $c$ is too expensive for slot $i$'') are fed back
\emph{downward} as domain reductions---exactly the kind of information
propagator cells consume.  In the lattice book's terms, reduced-cost fixing
is an additional propagator whose input is the LP solution and whose output
is a tightened domain.
\end{crossrefbox}

\subsection{Comparing Bound Quality}

Different bounding techniques give different quality bounds:

\begin{center}
\begin{tabular}{lcc}
\toprule
\textbf{Method} & \textbf{Bound quality} & \textbf{Computation cost} \\
\midrule
Propagation only & Weakest & Cheapest (milliseconds) \\
Lagrangian relaxation & Medium & Moderate (seconds) \\
LP relaxation & Tightest & Most expensive (per node) \\
\bottomrule
\end{tabular}
\end{center}

In practice, we use propagation at every B\&B node (cheap), LP relaxation at
the root and at nodes where propagation doesn't prune enough (moderate cost),
and Lagrangian relaxation as an alternative when the LP is too expensive.

\begin{exercisebox}
Implement the Lagrangian relaxation for the gear problem.  Run the
subgradient method for 100 iterations with a diminishing step size.
Compare the Lagrangian bound to the LP relaxation bound on the same
problem instance.  How close are they?
\end{exercisebox}

\begin{exercisebox}
Prove in Lean~4 that if the Lagrangian dual bound equals the incumbent's
penalty, then the incumbent is optimal.  Hint: use the weak Lagrangian
duality theorem.
\end{exercisebox}

% ════════════════════════════════════════════════════════════════════════════
%  PART III — HEURISTICS AND MULTI-OBJECTIVE
% ════════════════════════════════════════════════════════════════════════════
\part{Heuristics and Multi-Objective Optimisation}

% ────────────────────────────────────────────────────────────────────────────
\chapter{Local Search and Metaheuristics}
\label{ch:heuristics}
% ────────────────────────────────────────────────────────────────────────────

\epigraph{Perfect is the enemy of good---but good is the friend of perfect.}{}

Everything up to now has been exact: propagation is sound, LP bounds are
valid, B\&B is complete.  This chapter introduces \emph{heuristic} methods
that find good solutions fast, without guaranteeing optimality.  Their role
in the solver: provide a strong incumbent to warm-start B\&B, so bounding
can prune more aggressively from the start.

\section{Neighbourhood Structure}

\begin{defbox}[title={Neighbourhood for Gear Builds}]
Two builds are \emph{neighbours} if they differ in exactly one slot's stat
combo.  The \emph{neighbourhood} $N(b)$ of build $b$ is the set of all
builds reachable by changing one slot.  $|N(b)| \leq 16 \times 39 = 624$.
\end{defbox}

The neighbourhood defines the ``moves'' available to local search.  Its
size is a tradeoff: a larger neighbourhood (e.g., changing two slots at once,
giving $|N| \approx 16^2 \times 39^2 \approx 250{,}000$) explores more
options per step but takes longer per step.  The single-slot neighbourhood
is the standard choice for gear optimisation because it's large enough to
escape many local optima but small enough that evaluating all neighbours is
fast.

\begin{leanbox}[title={Lean 4 --- Neighbourhood Enumeration}]
\begin{minted}{lean}
/-- Generate all single-slot neighbours of a build. -/
def allNeighbours (b : Build) : List Build :=
  GearSlot.all.bind fun slot =>
    allCombos.filter (· ≠ b.gear slot) |>.map fun combo =>
      { b with gear := fun s => if s == slot then combo else b.gear s }

/-- A neighbour move: which slot changed and to what. -/
structure Move where
  slot  : GearSlot
  combo : StatCombo

/-- Generate all moves with their resulting penalties (for sorting). -/
def allMoves (b : Build) (targets : StatType → Nat)
    : List (Move × Nat) :=
  GearSlot.all.bind fun slot =>
    allCombos.filter (· ≠ b.gear slot) |>.map fun combo =>
      let b' := { b with gear := fun s =>
        if s == slot then combo else b.gear s }
      (⟨slot, combo⟩, totalPenalty b' targets wMiss wOver)
\end{minted}
\end{leanbox}

\section{Hill Climbing}

Start with a random feasible build.  Repeatedly move to the best neighbour
(lowest penalty).  Stop when no neighbour improves the current solution.

This is fast but easily trapped in local optima.  Restart with a new random
build and keep the best across restarts.

\begin{leanbox}[title={Lean 4 --- Hill Climbing with Restarts}]
\begin{minted}{lean}
/-- Hill climbing: greedily move to the best neighbour. -/
partial def hillClimb (current : Build) (currentPen : Nat)
    (targets : StatType → Nat) : Build × Nat :=
  let neighbours := allNeighbours current
  let (bestNeigh, bestPen) := neighbours.foldl
    (init := (current, currentPen)) fun (best, bestPen) neigh =>
      let pen := totalPenalty neigh targets wMiss wOver
      if pen < bestPen then (neigh, pen) else (best, bestPen)
  if bestPen < currentPen then hillClimb bestNeigh bestPen targets
  else (current, currentPen)  -- local optimum

/-- Restart hill climbing multiple times, keep the best. -/
def hillClimbRestarts (nRestarts : Nat) (rng : StdGen)
    (targets : StatType → Nat) : Build × Nat :=
  (List.range nRestarts).foldl (init := (defaultBuild, Nat.max, rng))
    fun (best, bestPen, rng) _ =>
      let (init, rng) := randomBuild rng
      let initPen := totalPenalty init targets wMiss wOver
      let (local, localPen) := hillClimb init initPen targets
      if localPen < bestPen then (local, localPen, rng)
      else (best, bestPen, rng)
  |>.fst
\end{minted}
\end{leanbox}

\section{Simulated Annealing}

\begin{defbox}[title={Simulated Annealing}]
Like hill climbing, but with a probability of accepting worse moves that
decreases over time:
\[
  P(\text{accept worse move with }\Delta > 0) = e^{-\Delta / T}
\]
where $T$ is the ``temperature,'' which decreases according to a cooling
schedule (e.g.\ $T_{k+1} = \alpha T_k$ with $\alpha \approx 0.995$).
\end{defbox}

\subsection{Choosing the Cooling Schedule}

The cooling schedule controls the exploration--exploitation tradeoff:

\begin{center}
\small
\begin{tabular}{lll}
\toprule
\textbf{Schedule} & \textbf{Formula} & \textbf{Character} \\
\midrule
Geometric & $T_{k+1} = \alpha T_k$ & Most common; $\alpha \in [0.99, 0.999]$ \\
Linear & $T_k = T_0 - k \cdot \delta$ & Simple but less effective \\
Logarithmic & $T_k = T_0 / \ln(k + 2)$ & Provably optimal but very slow \\
Adaptive & Adjust based on acceptance rate & Complex but robust \\
\bottomrule
\end{tabular}
\end{center}

For the gear problem, the geometric schedule with $\alpha = 0.997$ and
initial temperature $T_0$ set so that the initial acceptance probability
for a move with $\Delta = 100$ is about $0.8$ works well.  This gives
$T_0 = -100 / \ln(0.8) \approx 448$.

\begin{leanbox}[title={Lean 4 --- Simulated Annealing}]
\begin{minted}{lean}
/-- One step of simulated annealing. -/
def saStep (current : Build) (currentPen : Nat) (temp : Float)
    (rng : StdGen) (targets : StatType → Nat)
    : Build × Nat × StdGen :=
  -- Pick a random neighbour (random slot, random combo)
  let (slot, rng) := randomSlot rng
  let (combo, rng) := randomCombo rng (availableCombos slot)
  let neighbour := { current with gear := fun s =>
    if s == slot then combo else current.gear s }
  let neighPen := totalPenalty neighbour targets wMiss wOver
  if neighPen ≤ currentPen then
    (neighbour, neighPen, rng)  -- always accept improvements
  else
    let delta := (neighPen - currentPen : Float)
    let (r, rng) := randomFloat rng  -- r ∈ [0, 1)
    if r < Float.exp (-delta / temp) then
      (neighbour, neighPen, rng)  -- accept with probability
    else
      (current, currentPen, rng)  -- reject

/-- Full simulated annealing run. -/
def simulatedAnnealing (domains : GearSlot → Finset StatCombo)
    (targets : StatType → Nat) (rng : StdGen)
    (maxIters : Nat := 10000) : Build × Nat :=
  let (init, rng) := randomFeasibleBuild domains rng
  let initPen := totalPenalty init targets wMiss wOver
  let t0 : Float := 448.0    -- initial temperature
  let alpha : Float := 0.997  -- cooling rate
  (List.range maxIters).foldl
    (init := (init, initPen, init, initPen, t0, rng))
    fun (current, currentPen, best, bestPen, temp, rng) _ =>
      let (current', pen', rng') := saStep current currentPen temp rng targets
      let temp' := alpha * temp
      let (best', bestPen') :=
        if pen' < bestPen then (current', pen') else (best, bestPen)
      (current', pen', best', bestPen', temp', rng')
  |>.map fun (_, _, best, bestPen, _, _) => (best, bestPen)
\end{minted}
\end{leanbox}

\begin{histbox}
Simulated annealing was independently invented by Scott Kirkpatrick (IBM, 1983)
and Vlado \v{C}ern\'{y} (1985), inspired by the Metropolis--Hastings
algorithm from statistical physics.  The metallurgical analogy is precise:
annealing metal involves heating it (high temperature = accept bad moves)
and slowly cooling (low temperature = only accept improvements), allowing
atoms to find a low-energy crystalline structure rather than getting trapped
in an amorphous (locally optimal but globally suboptimal) state.
\end{histbox}

\section{Tabu Search}

\begin{defbox}[title={Tabu Search}]
Maintain a \emph{tabu list} of recently visited moves.  At each step, move
to the best non-tabu neighbour, even if it's worse than the current solution.
This prevents cycling and forces exploration past local optima.
\end{defbox}

\begin{leanbox}[title={Lean 4 --- Tabu Search}]
\begin{minted}{lean}
structure TabuState where
  current   : Build
  currentPen : Nat
  best      : Build
  bestPen   : Nat
  tabuList  : List (GearSlot × StatCombo)  -- recently changed
  tabuTenure : Nat                          -- how long moves stay tabu

def tabuStep (state : TabuState) (targets : StatType → Nat)
    : TabuState :=
  let candidates := allNeighbourMoves state.current
    |>.filter fun (slot, combo) =>
      -- Allow if not tabu, OR if it produces a new global best
      ¬ state.tabuList.contains (slot, combo)
  let (bestSlot, bestCombo, bestPen) := candidates.foldl
    (init := (.headgear, defaultCombo, Nat.max)) fun (bs, bc, bp) (slot, combo) =>
      let build := setBuildSlot state.current slot combo
      let pen := totalPenalty build targets wMiss wOver
      if pen < bp then (slot, combo, pen) else (bs, bc, bp)
  let newBuild := setBuildSlot state.current bestSlot bestCombo
  let newTabu := ((bestSlot, state.current.gear bestSlot) :: state.tabuList)
    |>.take state.tabuTenure  -- keep only recent entries
  { current := newBuild
    currentPen := bestPen
    best := if bestPen < state.bestPen then newBuild else state.best
    bestPen := min bestPen state.bestPen
    tabuList := newTabu
    tabuTenure := state.tabuTenure }
\end{minted}
\end{leanbox}

The tabu tenure (how long a move stays forbidden) is a key parameter.  Too
short and the search cycles; too long and it can't revisit promising regions.
A typical value for the gear problem: tenure $\approx \sqrt{16} \approx 4$
(the square root of the number of variables).

\section{Genetic Algorithms}

\begin{defbox}[title={Genetic Algorithm for Gear}]
\begin{itemize}[nosep]
  \item \textbf{Chromosome:} a build (assignment of combo to each slot).
  \item \textbf{Gene:} one slot's combo choice.
  \item \textbf{Crossover:} take armour slots from parent A, trinket slots
    from parent B.
  \item \textbf{Mutation:} randomly change one slot's combo.
  \item \textbf{Fitness:} $-\pen(b)$ (lower penalty = higher fitness).
  \item \textbf{Selection:} tournament selection (pick 3 random builds, keep
    the fittest).
\end{itemize}
\end{defbox}

\begin{leanbox}[title={Lean 4 --- Crossover and Mutation}]
\begin{minted}{lean}
/-- Uniform crossover: for each slot, pick from parent A or B. -/
def crossover (a b : Build) (rng : StdGen) : Build × StdGen :=
  GearSlot.all.foldl (init := (a, rng)) fun (child, rng) slot =>
    let (bit, rng) := randomBool rng
    let combo := if bit then a.gear slot else b.gear slot
    ({ child with gear := fun s => if s == slot then combo else child.gear s }, rng)

/-- Mutate: change one random slot to a random combo. -/
def mutate (b : Build) (rng : StdGen) (mutRate : Float := 0.1)
    : Build × StdGen :=
  let (r, rng) := randomFloat rng
  if r > mutRate then (b, rng)
  else
    let (slot, rng) := randomSlot rng
    let (combo, rng) := randomCombo rng allCombos
    ({ b with gear := fun s => if s == slot then combo else b.gear s }, rng)

/-- One generation of the GA. -/
def gaGeneration (pop : Array Build) (targets : StatType → Nat)
    (rng : StdGen) : Array Build × StdGen :=
  let fitnesses := pop.map (fun b => totalPenalty b targets wMiss wOver)
  (Array.range pop.size).foldl (init := (#[], rng)) fun (newPop, rng) _ =>
    -- Tournament selection (pick 3, keep best)
    let (parent1, rng) := tournamentSelect pop fitnesses 3 rng
    let (parent2, rng) := tournamentSelect pop fitnesses 3 rng
    let (child, rng) := crossover parent1 parent2 rng
    let (child, rng) := mutate child rng
    (newPop.push child, rng)
\end{minted}
\end{leanbox}

These heuristics are not verified for optimality.  But we \emph{can} verify
in Lean that their output is a \emph{feasible} build (satisfies all
constraints)---a weaker but still useful guarantee.

\begin{exercisebox}
Implement all three heuristics (SA, tabu, GA) for the gear problem.
Run each 10 times on the same target profile.  Compare: which finds
the best solution?  Which is most consistent?  Which is fastest?
\end{exercisebox}

\begin{exercisebox}
Prove in Lean~4: if \texttt{simulatedAnnealing} starts from a feasible build
(one where every slot's combo is in the initial domain), then it returns
a feasible build.  Hint: each step only changes one slot to a combo from
\texttt{availableCombos}.
\end{exercisebox}

% ────────────────────────────────────────────────────────────────────────────
\chapter{Multi-Objective Optimisation and Pareto Frontiers}
\label{ch:pareto}
% ────────────────────────────────────────────────────────────────────────────

\epigraph{You can't have it all---but you can know exactly what you're trading.}{}

\section{Pareto Optimality}

The gear problem is fundamentally multi-objective: maximise damage stats
(Power, Precision, Ferocity) \emph{and} survivability (Vitality, Toughness)
\emph{and} support (Healing Power, Concentration).

\begin{defbox}[title={Pareto Dominance}]
Build $a$ \emph{dominates} build $b$ (written $a \succ b$) if $a$ is at
least as good as $b$ in every stat and strictly better in at least one:
\[
  a \succ b \iff
  \big(\forall s.\; \statTotal(a, s) \geq \statTotal(b, s)\big)
  \;\wedge\;
  \big(\exists s.\; \statTotal(a, s) > \statTotal(b, s)\big).
\]
A build is \emph{Pareto optimal} if no feasible build dominates it.
\end{defbox}

\begin{crossrefbox}
The Pareto dominance relation is a \emph{strict partial order} (irreflexive,
transitive).  The set of all feasible builds, ordered by Pareto dominance,
forms a partially ordered set.  The Pareto frontier is the set of
\emph{maximal} elements---an \emph{antichain} (no element dominates another).

This connects directly to Chapter~2 of the lattice book: partial orders,
antichains, and maximal elements.
\end{crossrefbox}

\begin{propbox}[title={Pareto Dominance is a Strict Partial Order}]
The relation $\succ$ is irreflexive and transitive.
\end{propbox}

\begin{proof}
\emph{Irreflexive:} For any build $a$, we cannot have $a \succ a$:
the condition $\exists s.\; \statTotal(a, s) > \statTotal(a, s)$ is
always false.

\emph{Transitive:} Suppose $a \succ b$ and $b \succ c$.  Then for all $s$:
$\statTotal(a, s) \geq \statTotal(b, s) \geq \statTotal(c, s)$, so
$\statTotal(a, s) \geq \statTotal(c, s)$.  For strictness: since $a \succ b$,
there exists $s_1$ with $\statTotal(a, s_1) > \statTotal(b, s_1)$.  Since
$b \succ c$, $\statTotal(b, s_1) \geq \statTotal(c, s_1)$.  Therefore
$\statTotal(a, s_1) > \statTotal(c, s_1)$, giving the required strict
inequality.
\end{proof}

\begin{leanbox}[title={Lean 4 --- Pareto Dominance and Frontier}]
\begin{minted}{lean}
/-- Build a dominates build b. -/
def dominates (a b : Build) : Prop :=
  (∀ s : StatType, statTotal a s ≥ statTotal b s) ∧
  (∃ s : StatType, statTotal a s > statTotal b s)

/-- Dominance is irreflexive. -/
theorem dominates_irrefl (a : Build) : ¬ dominates a a := by
  intro ⟨_, ⟨s, hs⟩⟩
  exact Nat.lt_irrefl _ hs

/-- Dominance is transitive. -/
theorem dominates_trans {a b c : Build}
    (hab : dominates a b) (hbc : dominates b c) : dominates a c := by
  constructor
  · intro s
    exact Nat.le_trans (hbc.1 s) (hab.1 s) |>.symm
    -- actually: a ≥ b ≥ c
    sorry -- full proof uses Nat.le_trans
  · obtain ⟨s, hs⟩ := hab.2
    exact ⟨s, Nat.lt_of_lt_of_le hs (hbc.1 s) |>.symm⟩
    sorry -- exercise

/-- The Pareto frontier: all non-dominated builds. -/
def paretoFrontier (builds : List Build) : List Build :=
  builds.filter fun b =>
    ¬ builds.any fun a => dominates a b

/-- Decidable dominance for computation. -/
instance : Decidable (dominates a b) :=
  if h1 : StatType.all.all (fun s => statTotal a s ≥ statTotal b s) then
    if h2 : StatType.all.any (fun s => statTotal a s > statTotal b s) then
      isTrue ⟨by sorry, by sorry⟩
    else isFalse (by sorry)
  else isFalse (by sorry)
\end{minted}
\end{leanbox}

\subsection{Worked Example: Damage vs.\ Survivability}

Consider a simplified two-objective problem for a Warrior: maximise
\emph{effective power} (a function of Power, Precision, Ferocity) and
\emph{effective health} (a function of Vitality, Toughness).

\begin{center}
\begin{tikzpicture}[scale=0.8]
  \draw[-{Stealth}, thick] (0, 0) -- (8, 0) node[right] {Effective Power};
  \draw[-{Stealth}, thick] (0, 0) -- (0, 6) node[above] {Effective Health};
  % Pareto frontier (stylised)
  \fill[defbg, opacity=0.5] (1,5) -- (2,4.5) -- (3.5,3.8) -- (5,2.5) -- (6.5,1.5) -- (7,0.5) -- (8,0) -- (0,0) -- (0,6) -- cycle;
  \draw[midblue, very thick] (1,5) -- (2,4.5) -- (3.5,3.8) -- (5,2.5) -- (6.5,1.5) -- (7,0.5);
  % Points
  \fill[accentgreen] (1,5) circle (3pt) node[above right, font=\scriptsize] {Full Minstrel's};
  \fill[accentgreen] (2,4.5) circle (3pt) node[above right, font=\scriptsize] {Minstrel/Marauder mix};
  \fill[accentgreen] (3.5,3.8) circle (3pt) node[right, font=\scriptsize] {Celestial};
  \fill[accentgreen] (5,2.5) circle (3pt) node[right, font=\scriptsize] {Marauder};
  \fill[accentgreen] (6.5,1.5) circle (3pt) node[right, font=\scriptsize] {Berserker/Marauder};
  \fill[accentgreen] (7,0.5) circle (3pt) node[right, font=\scriptsize] {Full Berserker's};
  % Dominated point
  \fill[accentred] (4,2) circle (3pt) node[right, font=\scriptsize] {Dominated};
  \draw[accentred, dashed] (4,2) -- (5,2.5);
  \node[font=\small, text=gray] at (4, -0.7) {Pareto frontier in \textcolor{midblue}{blue}; dominated point in \textcolor{accentred}{red}};
\end{tikzpicture}
\end{center}

The Pareto frontier traces the tradeoff: full Berserker's maximises damage
but has minimal health; full Minstrel's maximises health but deals little
damage.  The ``knee'' (where the frontier curves most sharply) often
corresponds to a well-balanced build---in this case, Marauder or Celestial.

A point like the dominated one (red) is strictly worse than Marauder in
\emph{both} dimensions: more damage \emph{and} more health are available.
The solver would never recommend a dominated build.

\section{Scalarisation}

Different penalty weights $w_{\text{miss}}$ and $w_{\text{over}}$, or
different target profiles, trace different points on the Pareto frontier.
By sweeping the weights, we can approximate the full frontier.

\begin{defbox}[title={Weighted Scalarisation}]
Given weight vector $w = (w_1, \ldots, w_k) \geq 0$ with $\sum w_i = 1$:
\[
  \min_b \sum_{i=1}^{k} w_i \cdot f_i(b)
\]
where each $f_i$ is a single-objective penalty.  Different choices of $w$
yield different Pareto-optimal solutions.
\end{defbox}

\begin{warningbox}
Weighted scalarisation can only find points on the \emph{convex} part of the
Pareto frontier.  Non-convex portions (concavities in the frontier) are
invisible to linear scalarisation.  The $\varepsilon$-constraint method
handles this.
\end{warningbox}

\section{The $\varepsilon$-Constraint Method}

Fix all objectives except one as constraints: ``optimise Power subject to
Vitality $\geq V_{\min}$.''  Sweep $V_{\min}$ to trace the frontier.

\begin{leanbox}[title={Lean 4 --- $\varepsilon$-Constraint Sweep}]
\begin{minted}{lean}
/-- Trace the Pareto frontier by sweeping an epsilon constraint. -/
def epsilonConstraintSweep
    (problem : GearProblem)
    (primaryStat : StatType)      -- optimise this
    (constrainedStat : StatType)  -- constrain this
    (minVal maxVal step : Nat)    -- sweep range
    : List (Nat × Nat × Build) :=  -- (primary, constrained, build)
  (List.range ((maxVal - minVal) / step + 1)).filterMap fun i =>
    let threshold := minVal + i * step
    let constrained := { problem with
      targets := fun s =>
        if s == constrainedStat then threshold
        else problem.targets s }
    match solve constrained with
    | .optimal build pen => some (statTotal build primaryStat,
                                  statTotal build constrainedStat, build)
    | .infeasible => none
\end{minted}
\end{leanbox}

\begin{exercisebox}
For a Warrior with base stats Power~1000, Precision~1000, Vitality~1000,
Toughness~1000: use the $\varepsilon$-constraint method to trace the
Pareto frontier between effective power (Power + Precision + Ferocity)
and effective health (Vitality $\times$ (1 + Toughness / 1000)).  Plot
the resulting frontier.
\end{exercisebox}

% ────────────────────────────────────────────────────────────────────────────
\chapter{Infusion Optimisation}
\label{ch:infusion}
% ────────────────────────────────────────────────────────────────────────────

\epigraph{The last mile is the easiest---if you plan it right.}{}

Once gear slots are fixed, the remaining problem is: distribute 18 infusion
slots (each $+5$ to one stat) to minimise the remaining penalty.  This is a
small, clean subproblem.

\section{The Resource Allocation Formulation}

\begin{defbox}[title={Infusion Subproblem}]
Given fixed gear contributions $G(s)$ for each stat $s$, and targets $T(s)$:
\[
  \min_{n_1, \ldots, n_k} \sum_{s} \pen\!\big(T(s),\; G(s) + 5 n_s\big)
  \quad\text{s.t.}\quad
  \sum_s n_s = 18, \quad n_s \geq 0, \quad n_s \in \N.
\]
\end{defbox}

\section{Dynamic Programming Solution}

\begin{defbox}[title={Infusion DP}]
Define $\texttt{dp}[i][r]$ = minimum penalty using stats $1, \ldots, i$
with $r$ infusions remaining.  The recurrence:
\[
  \texttt{dp}[i][r] = \min_{0 \leq n_i \leq r} \;
    \pen(T(s_i), G(s_i) + 5 n_i) + \texttt{dp}[i-1][r - n_i].
\]
Base case: $\texttt{dp}[0][r] = 0$ for all $r$.
\end{defbox}

\begin{leanbox}[title={Lean 4 --- Infusion DP}]
\begin{minted}{lean}
/-- Optimally distribute infusions to minimise remaining penalty. -/
def infusionDP (gearStats : StatType → Nat)
    (targets : StatType → Nat) (budget : Nat := 18)
    : StatType → Nat :=
  let stats := StatType.all.toArray
  -- dp[i][r] = (minPenalty, allocation for stats 0..i)
  let dp := Array.mkArray (stats.size + 1) (Array.mkArray (budget + 1) (0, #[]))
  let dp := stats.foldl (init := dp) fun dp i stat =>
    Array.mkArray (budget + 1) (0, #[]) |>.mapIdx fun r _ =>
      let best := (List.range (r.val + 1)).foldl (init := (Nat.max, #[])) fun acc ni =>
        let penHere := penalty (targets stat) (gearStats stat + 5 * ni) wMiss wOver
        let (penRest, allocRest) := dp[i]![r.val - ni]!
        let total := penHere + penRest
        if total < acc.1 then (total, allocRest.push ni) else acc
      best
  -- Extract the allocation from dp[k][budget]
  let (_, alloc) := dp[stats.size]![budget]!
  fun stat => alloc[stats.indexOf stat]!
\end{minted}
\end{leanbox}

\begin{thmbox}[title={Optimality of Infusion DP}]
The DP computes the optimal infusion distribution: for any fixed gear
assignment, the allocation returned minimises the total penalty.
\end{thmbox}

\begin{proof}
By induction on the number of stat types.  The base case (0 stats) is
trivial.  The inductive step: the DP considers all possible allocations
$n_i$ for stat $i$ and recursively computes the optimal allocation for
the remaining stats with the reduced budget.  By the inductive hypothesis,
each subproblem is solved optimally, so the minimum over $n_i$ choices
yields the global optimum.
\end{proof}

The DP has $O(k \cdot B)$ states (where $k = 9$ stat types and $B = 18$
infusion budget), and $O(B)$ transitions per state, for a total of
$O(k \cdot B^2) = O(9 \cdot 324) = O(2916)$ operations.  This is negligible
and can be called at every leaf of the B\&B tree without performance concern.

% ════════════════════════════════════════════════════════════════════════════
%  PART IV — THE CAPSTONE SOLVER
% ════════════════════════════════════════════════════════════════════════════
\part{The Capstone Solver}

% ────────────────────────────────────────────────────────────────────────────
\chapter{Assembly: The Complete Solver}
\label{ch:assembly}
% ────────────────────────────────────────────────────────────────────────────

\epigraph{The whole is greater than the sum of its parts.}{}

We now wire together every component into a single solver.

\section{The Architecture}

\begin{center}
\begin{tikzpicture}[
  box/.style={draw=midblue, fill=defbg, rounded corners=3pt,
    minimum width=8cm, minimum height=0.9cm, font=\small, align=center},
  smallbox/.style={draw=accentgreen!70, fill=thmbg, rounded corners=3pt,
    minimum width=3.5cm, minimum height=0.7cm, font=\scriptsize, align=center},
  arr/.style={-{Stealth[length=5pt]}, thick, midblue},
  note/.style={font=\scriptsize\itshape, text=gray}
]
\node[box] (warm) at (0, 6) {\textbf{1. Warm Start} (SA/GA, Ch.~\ref{ch:heuristics})};
\node[box] (prop) at (0, 4.5) {\textbf{2. Initial Propagation} (Ch.~\ref{ch:stronger})};
\node[box] (bb)   at (0, 3)   {\textbf{3. Branch \& Bound} (Ch.~\ref{ch:bb})};
\node[smallbox] (lp)   at (-3, 1.5)  {LP Bound\\(Ch.~\ref{ch:lp})};
\node[smallbox] (prop2) at (0, 1.5) {Propagation\\at each node};
\node[smallbox] (inf)  at (3, 1.5)  {Infusion DP\\(Ch.~\ref{ch:infusion})};
\node[box] (out) at (0, 0) {\textbf{4. Result:} Optimal build + proof, or infeasibility};

\draw[arr] (warm) -- (prop);
\draw[arr] (prop) -- (bb);
\draw[arr] (bb) -- (lp);
\draw[arr] (bb) -- (prop2);
\draw[arr] (bb) -- (inf);
\draw[arr] (lp) -- (out);
\draw[arr] (prop2) -- (out);
\draw[arr] (inf) -- (out);
\end{tikzpicture}
\end{center}

\section{The Solver Loop}

\begin{leanbox}[title={Lean 4 --- The Complete Solver}]
\begin{minted}{lean}
inductive SolverResult where
  | optimal    : Build → Nat → SolverResult   -- build + penalty
  | infeasible : SolverResult

def solve (problem : GearProblem) : SolverResult :=
  -- Phase 1: Initial propagation
  let domains := propagateToFixpoint problem.initialDomains
  if anyDomainEmpty domains then .infeasible
  else
  -- Phase 2: Warm start with simulated annealing
  let warmBuild := simulatedAnnealing domains problem.targets
  let incumbent := (warmBuild, totalPenalty warmBuild problem.targets wMiss wOver)
  -- Phase 3: Branch and bound
  let result := branchAndBound domains problem.targets incumbent
  match result with
  | none => .infeasible
  | some (build, pen) => .optimal build pen

/-- Branch and bound with propagation + LP bounding at each node. -/
partial def branchAndBound
    (domains : GearSlot → Finset StatCombo)
    (targets : StatType → Nat)
    (incumbent : Build × Nat)
    : Option (Build × Nat) :=
  -- Propagate at this node
  let domains := propagateToFixpoint domains
  if anyDomainEmpty domains then some incumbent
  -- All singleton? → Evaluate (with infusion DP)
  else if allSingleton domains then
    let build := extractBuild domains
    let infAlloc := infusionDP (gearStats build) targets
    let fullBuild := { build with infusions := infAlloc }
    let pen := totalPenalty fullBuild targets wMiss wOver
    if pen < incumbent.2 then some (fullBuild, pen)
    else some incumbent
  -- LP bound check
  else
    let lpBound := solveLPRelaxation domains targets
    if lpBound ≥ incumbent.2 then some incumbent  -- prune
    else
      -- Branch: fail-first variable ordering
      let slot := smallestDomainSlot domains
      (domains slot).fold incumbent fun best combo =>
        let domains' := fixSlot domains slot combo
        match branchAndBound domains' targets best with
        | none => best
        | some result => result
\end{minted}
\end{leanbox}

% ────────────────────────────────────────────────────────────────────────────
\chapter{The Correctness Theorem}
\label{ch:correctness}
% ────────────────────────────────────────────────────────────────────────────

\epigraph{Trust, but verify.}{}

The verified core.  We state and prove the main theorem: the solver is
\emph{sound} (if it says optimal, it is; if it says infeasible, no build
works).

\section{The Main Theorem}

\begin{thmbox}[title={Solver Soundness}]
\begin{minted}{lean}
theorem solver_sound (problem : GearProblem) :
    match solve problem with
    | .optimal build pen =>
        build.feasible problem ∧
        pen = totalPenalty build problem.targets wMiss wOver ∧
        ∀ b, b.feasible problem →
          totalPenalty b problem.targets wMiss wOver ≥ pen
    | .infeasible =>
        ∀ b : Build, ¬ b.feasible problem
\end{minted}
\end{thmbox}

\section{Proof Structure}

The proof composes four lemmas.  We present each with its Lean formalisation
and proof sketch.

\subsection{Propagation Soundness}

\begin{lembox}[title={Propagation Preserves Optimal Solutions}]
If a combo $c$ is removed from slot $i$'s domain during propagation, then
no optimal solution uses $c$ on slot $i$.
\end{lembox}

\begin{proof}
By the monotonicity of each propagator (proven in Chapter~\ref{ch:stronger}).
The \texttt{LinearGeq} propagator removes $c$ from $D_i$ only when
\[
  \texttt{contribution}(i, c, s) + \sum_{j \neq i} \maxContrib(j, s) + 90 < T(s)
\]
for some stat $s$.  This means no assignment using $c$ on slot $i$ can meet
target $T(s)$, even with all other slots maximised and all infusions devoted
to $s$.  Any build using $c$ on slot $i$ violates the target, so its penalty
is strictly positive.  If a zero-penalty build exists, it doesn't use $c$ on
slot $i$.  If no zero-penalty build exists, the propagator is still sound:
it may remove suboptimal choices, but never the optimal one (since the optimal
choice must be at least as good as $c$ on the pruned stat).
\end{proof}

\begin{leanbox}[title={Lean 4 --- Propagation Soundness}]
\begin{minted}{lean}
/-- If LinearGeq removes combo c from slot i, then any build using c
    on slot i has penalty > 0 on some stat. -/
theorem linearGeq_sound
    (domains : GearSlot → Finset StatCombo)
    (targets : StatType → Nat)
    (slot : GearSlot) (combo : StatCombo)
    (h_removed : combo ∉ linearGeqPropagate domains targets slot)
    (h_was_in : combo ∈ domains slot)
    (b : Build) (h_uses : b.gear slot = combo) :
    ∃ s : StatType, statTotal b s + 90 < targets s := by
  -- h_removed gives us: ∃ stat, contribution + othersMax + 90 < target
  -- Since statTotal b s ≤ contribution(slot, combo, s) + othersMax(s),
  -- the result follows
  sorry -- full proof uses the definition of linearGeqPropagate

/-- Propagation to fixpoint preserves all optimal solutions. -/
theorem propagation_preserves_optimal
    (domains₀ : GearSlot → Finset StatCombo)
    (targets : StatType → Nat)
    (b_opt : Build)
    (h_opt : isOptimal problem b_opt)
    (h_in : ∀ s, b_opt.gear s ∈ domains₀ s) :
    let domains' := propagateToFixpoint domains₀
    ∀ s, b_opt.gear s ∈ domains' s := by
  -- By induction on the number of propagation steps
  -- At each step, the combo removed can't be part of b_opt
  -- (by linearGeq_sound and h_opt)
  sorry
\end{minted}
\end{leanbox}

\subsection{LP Bound Validity}

\begin{lembox}[title={LP Relaxation is a Valid Lower Bound}]
For any node $n$ in the B\&B tree:
$\LP(n) \leq \min\{\pen(b) \mid b \text{ is feasible in subtree } n\}$.
\end{lembox}

\begin{proof}
The LP relaxation enlarges the feasible set (from integer to fractional
assignments).  The minimum over a larger set is $\leq$ the minimum over
the original set.

Formally: let $F_{\IP}(n) = \{b \in \text{Builds} \mid b.gear(i) \in D_i(n)\;\forall i\}$
be the integer-feasible builds at node $n$, and $F_{\LP}(n)$ be the LP-feasible
fractional assignments.  Since every integer assignment is also a valid
fractional assignment, $F_{\IP}(n) \subseteq F_{\LP}(n)$.  Therefore:
\[
  \LP(n) = \min_{x \in F_{\LP}(n)} \pen(x) \leq \min_{b \in F_{\IP}(n)} \pen(b).
\]
\end{proof}

\begin{leanbox}[title={Lean 4 --- LP Bound Validity}]
\begin{minted}{lean}
/-- The LP relaxation value is a lower bound on any integer solution. -/
theorem lp_bound_valid
    (node : BBNode)
    (lpResult : LPResult)
    (h_lp : lpResult = solveLPRelaxation node.domains node.targets)
    (b : Build)
    (h_feasible : ∀ s, b.gear s ∈ node.domains s) :
    lpResult.objectiveValue ≤ (totalPenalty b node.targets wMiss wOver).toFloat := by
  -- The build b corresponds to a 0-1 assignment, which is feasible for the LP
  -- The LP optimum ≤ any feasible LP solution
  sorry
\end{minted}
\end{leanbox}

\subsection{B\&B Invariant}

\begin{lembox}[title={B\&B Invariant Preservation}]
At every point during the search: every feasible build either (a) has been
evaluated and is no better than the incumbent, or (b) lives in an unexplored
node whose LP bound is $\leq$ the incumbent penalty.
\end{lembox}

\begin{proof}
By induction on the search steps.

\emph{Base case:} before any branching, the entire search space is a single
unexplored node.  All feasible builds live there.

\emph{Inductive step:} at each step, one of three things happens:
\begin{enumerate}[nosep]
  \item \textbf{Branch:} an unexplored node is split into children.
    Every feasible build in the parent lives in exactly one child.
    The invariant is preserved because the children partition the parent.
  \item \textbf{Prune:} a node is discarded because $\LP(n) \geq
    \texttt{incumbent}$.  By LP bound validity, every integer solution
    in this subtree has penalty $\geq \LP(n) \geq \texttt{incumbent}$.
    These builds are accounted for.
  \item \textbf{Evaluate:} a leaf node (all singletons) is evaluated.
    If its penalty $<$ incumbent, it becomes the new incumbent.  Otherwise,
    it's no better.
\end{enumerate}
In all cases, every feasible build is either evaluated or in an accounted
node.
\end{proof}

\begin{leanbox}[title={Lean 4 --- B\&B Invariant}]
\begin{minted}{lean}
/-- The B&B invariant: incumbent is feasible and dominates all pruned nodes. -/
structure BBInvariant (state : BBState) where
  incumbent_feasible : state.incumbent.feasible state.problem
  pruned_bounded : ∀ node ∈ state.pruned,
    node.lpBound ≥ state.incumbent.penalty
  partition : ∀ b : Build, b.feasible state.problem →
    (totalPenalty b ≥ state.incumbent.penalty) ∨
    (∃ node ∈ state.queue, b.inSubtree node)

/-- Each step of B&B preserves the invariant. -/
theorem bb_step_preserves (state : BBState)
    (h_inv : BBInvariant state)
    (state' : BBState)
    (h_step : state' = bbStep state) :
    BBInvariant state' := by
  sorry -- case split on branch/prune/evaluate
\end{minted}
\end{leanbox}

\subsection{Infusion DP Optimality}

\begin{lembox}[title={Infusion DP Optimality}]
For any fixed gear assignment, the infusion DP returns the allocation
that minimises the total penalty.
\end{lembox}

\begin{proof}
By induction on the number of stat types $k$.

\emph{Base case ($k = 0$):} no stats to allocate infusions to, penalty is
determined entirely by gear.  The DP returns the empty allocation with the
correct penalty.

\emph{Inductive step:} assume the DP is optimal for stats $1, \ldots, k-1$.
For stat $k$, the DP tries every possible number of infusions
$n_k \in \{0, \ldots, \text{remaining}\}$ and recursively computes the
optimal allocation for the remaining stats with the reduced budget.  By the
inductive hypothesis, each recursive call returns the optimal sub-allocation.
Therefore, the minimum over all choices of $n_k$ yields the global optimum.

\emph{Why this works:} the penalty function is separable across stats
($\pen = \sum_s \pen_s$), so the optimal total allocation decomposes into
per-stat choices given a budget.  This is exactly the optimal substructure
property that dynamic programming exploits.
\end{proof}

\begin{leanbox}[title={Lean 4 --- Infusion DP Optimality Proof}]
\begin{minted}{lean}
/-- The DP allocation minimises penalty among all valid allocations. -/
theorem infusion_dp_optimal
    (gearStats : StatType → Nat) (targets : StatType → Nat)
    (budget : Nat) :
    let dpAlloc := infusionDP gearStats targets budget
    ∀ alloc : StatType → Nat,
      (StatType.all.map alloc).sum = budget →
      totalInfusionPenalty gearStats targets dpAlloc ≤
      totalInfusionPenalty gearStats targets alloc := by
  -- By induction on StatType.all.length (= 9)
  -- using the recurrence relation of the DP table
  sorry
\end{minted}
\end{leanbox}

\subsection{Composing the Proof}

When the B\&B search completes (queue empty), the invariant implies:
\begin{enumerate}[nosep]
  \item The incumbent is feasible (it was constructed from a valid assignment
    and optimal infusions, by infusion DP optimality).
  \item Every other feasible build either was evaluated (and is no better)
    or lives in a pruned subtree (whose LP bound $\geq$ incumbent, by LP
    bound validity).
  \item By propagation soundness, all pruned combos were provably suboptimal,
    so the search space explored contains all optimal solutions.
  \item Therefore, the incumbent is optimal.
\end{enumerate}

The infeasibility case: if initial propagation empties a domain, no
assignment exists.  By propagation soundness, every pruned combo was
provably useless, so the empty domain means \emph{no} feasible build exists.

\begin{leanbox}[title={Lean 4 --- Composing the Full Proof}]
\begin{minted}{lean}
/-- The complete solver soundness theorem. -/
theorem solver_sound_proof (problem : GearProblem) :
    match solve problem with
    | .optimal build pen =>
        build.feasible problem ∧
        pen = totalPenalty build problem.targets wMiss wOver ∧
        ∀ b, b.feasible problem →
          totalPenalty b problem.targets wMiss wOver ≥ pen
    | .infeasible =>
        ∀ b : Build, ¬ b.feasible problem := by
  unfold solve
  -- Case 1: initial propagation finds empty domain
  · -- by propagation_preserves_optimal, no feasible build existed
    intro b hfeas
    exact absurd (propagation_preserves_optimal ..) (domain_empty ..)
  -- Case 2: B&B returns a solution
  · constructor
    -- feasibility: build was extracted from valid domains
    · exact extractBuild_feasible ..
    constructor
    -- penalty matches: by definition of totalPenalty
    · rfl
    -- optimality: by B&B invariant + LP bound validity
    · intro b hfeas
      exact bb_invariant_at_termination .. |>.partition b hfeas |>.elim
        (fun h => h)
        (fun ⟨node, _, _⟩ => absurd .. (queue_empty ..))
\end{minted}
\end{leanbox}

This is the central theorem of the book.  The \texttt{sorry}s in the
component lemmas indicate where the reader should fill in the full
proofs as exercises---each one is a substantial but tractable formalisation
task that applies techniques from the lattice book.

% ────────────────────────────────────────────────────────────────────────────
\chapter{Performance Engineering}
\label{ch:perf}
% ────────────────────────────────────────────────────────────────────────────

\epigraph{An elegant algorithm with a slow implementation\\
is just a slow algorithm.}{}

\section{Domain Representation}

Replace \texttt{Finset StatCombo} with \texttt{BitVec 40} (one bit per
combo).  Set intersection becomes bitwise AND, union becomes OR, cardinality
is \texttt{popcount}.  This makes propagation dramatically faster.

\begin{leanbox}[title={Lean 4 --- BitVec Domain}]
\begin{minted}{lean}
/-- A domain is a 40-bit vector: bit i = 1 iff combo i is still possible. -/
abbrev Domain := BitVec 40

def Domain.intersect (a b : Domain) : Domain := a &&& b
def Domain.isEmpty (d : Domain) : Bool := d == 0
def Domain.card (d : Domain) : Nat := d.popcount
def Domain.remove (d : Domain) (i : Fin 40) : Domain :=
  d &&& ~~~(1 <<< i.val)
\end{minted}
\end{leanbox}

The performance difference is substantial.  \texttt{Finset} operations
involve linked-list traversals; \texttt{BitVec} operations are single
machine instructions:

\begin{center}
\small
\begin{tabular}{lcc}
\toprule
\textbf{Operation} & \textbf{Finset} & \textbf{BitVec 40} \\
\midrule
Intersection ($D_1 \cap D_2$) & $O(|D_1| + |D_2|)$ & $O(1)$ (AND) \\
Is empty? & $O(1)$ & $O(1)$ (compare to 0) \\
Cardinality & $O(|D|)$ & $O(1)$ (popcount) \\
Remove element & $O(|D|)$ & $O(1)$ (AND NOT mask) \\
Contains? & $O(|D|)$ & $O(1)$ (AND mask, compare) \\
\bottomrule
\end{tabular}
\end{center}

Since propagation runs at every B\&B node, and each propagation round does
$O(\text{slots} \times \text{stats} \times \text{domain size})$ operations,
the BitVec representation gives a constant-factor speedup of roughly $10$--$40\times$
on domain operations.

\begin{leanbox}[title={Lean 4 --- Precomputed Contribution Tables}]
\begin{minted}{lean}
/-- Precompute contribution of each combo to each stat on each slot.
    Access: contribTable[slot][combo][stat] -/
def buildContribTable : Array (Array (Array Nat)) :=
  GearSlot.all.toArray.map fun slot =>
    allCombos.toArray.map fun combo =>
      StatType.all.toArray.map fun stat =>
        contribution slot combo stat

/-- Precompute max contribution per stat per domain.
    Avoids recomputing during propagation. -/
structure DomainStats where
  maxPerStat  : Array Nat   -- maxPerStat[stat] = max contribution in domain
  minPerStat  : Array Nat   -- minPerStat[stat] = min contribution in domain
  sumMaxPower : Nat         -- cached sum for quick feasibility check
\end{minted}
\end{leanbox}

\section{Incremental Propagation}

Instead of re-running all propagators from scratch at each B\&B node,
maintain a queue of ``dirty'' cells (slots whose domains changed since
the last propagation round).  Only re-fire propagators that read from
dirty cells.

\begin{leanbox}[title={Lean 4 --- Incremental Propagation Queue}]
\begin{minted}{lean}
structure PropState where
  domains    : Array Domain           -- one per slot (16 entries)
  dirty      : Array Bool             -- which slots changed?
  domStats   : Array DomainStats      -- cached per-slot stats
  feasible   : Bool                   -- still feasible?

/-- Propagate incrementally: only process dirty slots. -/
def propagateIncremental (state : PropState)
    (targets : StatType → Nat) : PropState :=
  let dirtySlots := state.dirty.toList.enum.filter (·.2) |>.map (·.1)
  if dirtySlots.isEmpty then state  -- fixpoint reached
  else
    let state' := dirtySlots.foldl (init := { state with dirty := state.dirty.map (fun _ => false) })
      fun st slotIdx =>
        -- Re-run LinearGeq for all slots that might be affected
        GearSlot.all.foldl (init := st) fun st2 slot =>
          let newDomain := linearGeqFilter st2.domains slot targets
          if newDomain != st2.domains[slot.toIdx]! then
            { st2 with
              domains := st2.domains.set! slot.toIdx newDomain
              dirty := st2.dirty.set! slot.toIdx true }
          else st2
    propagateIncremental state' targets  -- recurse until clean
\end{minted}
\end{leanbox}

The key insight: when branching (fixing a slot to a combo), only \emph{that}
slot's domain changes.  Incremental propagation avoids redundantly re-checking
slots whose domains haven't changed.

\section{Undo Stacks for Backtracking}

Rather than copying the entire domain array at each B\&B node, maintain an
\emph{undo stack}: record which domains changed and what their old values were.
When backtracking, pop the stack and restore.

\begin{leanbox}[title={Lean 4 --- Undo Stack}]
\begin{minted}{lean}
structure UndoEntry where
  slot     : Fin 16
  oldDomain : Domain

/-- Save the current domain before a branch. -/
def saveState (domains : Array Domain) (slot : Fin 16)
    : UndoEntry :=
  ⟨slot, domains[slot]!⟩

/-- Restore domain after backtracking. -/
def restoreState (domains : Array Domain) (entry : UndoEntry)
    : Array Domain :=
  domains.set! entry.slot entry.oldDomain
\end{minted}
\end{leanbox}

This eliminates the $O(16 \times 40)$ cost of copying domain arrays at each
node, replacing it with $O(1)$ per save and restore.

\section{Lean 4 Performance Annotations}

\begin{leanbox}[title={Lean 4 --- Compiler Hints}]
\begin{minted}{lean}
-- Force inlining of hot inner loops
@[inline] def contribution' (slot : GearSlot) (comboIdx : Fin 40)
    (statIdx : Fin 9) : Nat :=
  contribTable[slot.toIdx]![comboIdx]![statIdx]!

-- Specialise generic functions for our concrete types
@[specialize] def propagateWith [BEq D] [Domain D]
    (domains : Array D) (props : Array (Propagator D)) : Array D :=
  sorry

-- Use unboxed representations where possible
structure PackedBuild where
  slots     : BitVec 240  -- 16 slots × 6 bits each (enough for 40 combos)
  runeIdx   : UInt8
  infusions : BitVec 72   -- 9 stats × 8 bits each (max 18 per stat)
\end{minted}
\end{leanbox}

\section{Benchmarks}

Run the full solver on realistic GW2 data: all 40 combos, all slot types,
several target profiles.  Compare:

\begin{center}
\small
\begin{tabular}{lccc}
\toprule
\textbf{Method} & \textbf{Nodes explored} & \textbf{Time} & \textbf{Optimal?} \\
\midrule
Brute force & $40^{16}$ (infeasible) & --- & Yes \\
Propagation + enumerate & $\sim 10^{12}$ & hours & Yes \\
Propagation + B\&B (no LP) & $\sim 10^{6}$ & seconds & Yes \\
Full solver (prop + B\&B + LP) & $\sim 10^{3}$ & milliseconds & Yes \\
SA warm-start + full solver & $\sim 10^{2}$ & milliseconds & Yes \\
SA alone (heuristic) & $10^{4}$ iterations & milliseconds & No \\
\bottomrule
\end{tabular}
\end{center}

The numbers are illustrative---actual performance depends on the target
profile and how constrained the problem is.  Tight targets (close to the
maximum achievable) produce more pruning and faster solves.  Loose targets
leave more feasible space to explore.

\begin{exercisebox}
Implement the \texttt{BitVec 40} domain representation and the precomputed
contribution table.  Run the strengthened propagator network on a standard
test case (Power $\geq 2500$, Precision $\geq 2000$, Ferocity $\geq 1200$
for heavy armour).  Measure the time difference vs.\ the \texttt{Finset}
representation.
\end{exercisebox}

\begin{exercisebox}
Implement the undo stack and incremental propagation.  Compare the number
of propagation steps per B\&B node with and without incrementality on a
10-slot problem (6 armour + 4 trinkets).
\end{exercisebox}

% ────────────────────────────────────────────────────────────────────────────
\chapter{Extensions and Open Problems}
\label{ch:extensions}
% ────────────────────────────────────────────────────────────────────────────

\epigraph{Every solved problem opens the door to a harder one.}{}

\section{Food and Utility Buffs}

Food and utility consumables provide flat stat bonuses.  A player selects
one food and one utility, each from a finite menu.  Mechanically, these are
two additional ``slots'' in the optimisation problem:

\begin{leanbox}[title={Lean 4 --- Modelling Food and Utility}]
\begin{minted}{lean}
structure FoodItem where
  name  : String
  stats : StatType → Nat  -- flat bonuses (e.g., +100 Power, +70 Ferocity)

structure UtilityItem where
  name  : String
  stats : StatType → Nat  -- flat bonuses (e.g., +100 Precision)

/-- Extended build includes food and utility. -/
structure ExtBuild extends Build where
  food    : FoodItem
  utility : UtilityItem

/-- Extended stat total includes food/utility contributions. -/
def extStatTotal (b : ExtBuild) (s : StatType) : Nat :=
  statTotal b.toBuild s + b.food.stats s + b.utility.stats s
\end{minted}
\end{leanbox}

The solver handles food and utility naturally.  They become two additional
variables with finite domains (the list of available foods and utilities).
Propagation extends: the \texttt{LinearGeq} propagator now includes the
maximum possible food/utility contribution when checking stat floors.  Since
food and utility have no coupling constraints with gear (unlike the rune
constraint on armour), they don't complicate the structure significantly.

\begin{intuitionbox}
Because food and utility affect \emph{every} stat simultaneously, fixing them
early in the search tree can dramatically improve propagation effectiveness.
A good heuristic: branch on food first, then utility, then gear slots.
This gives propagation a concrete base of flat bonuses to reason about.
\end{intuitionbox}

\begin{exercisebox}
Extend the solver to include food and utility.  How does adding these two
variables change the size of the search tree?  If there are 30 foods and
20 utilities, the tree grows by a factor of $30 \times 20 = 600$---but
how much does the improved propagation (tighter stat floors) compensate?
\end{exercisebox}

\section{Trait Interactions}

Some profession traits convert one stat into another.  For example, the
Guardian trait \emph{Retribution} might grant Power equal to a percentage
of Toughness.  This creates a \emph{nonlinear} dependency:

\begin{leanbox}[title={Lean 4 --- Trait-Modified Stat Computation}]
\begin{minted}{lean}
structure TraitConversion where
  sourceStat : StatType
  targetStat : StatType
  percentage : Float   -- e.g., 0.07 for 7%

/-- Stat total with trait conversions applied. -/
def traitStatTotal (b : Build) (conversions : List TraitConversion)
    (s : StatType) : Float :=
  let base := (statTotal b s).toFloat
  let bonus := conversions.foldl (init := 0.0) fun acc conv =>
    if conv.targetStat == s then
      conv.percentage * (statTotal b conv.sourceStat).toFloat
    else acc
  base + bonus
\end{minted}
\end{leanbox}

The difficulty: with trait conversions, adding Toughness gear increases Power
(through the trait), which means the stat computation is no longer a simple
sum over independent slot contributions.  The LP relaxation assumes linearity
and becomes invalid.

\paragraph{Approach 1: Linearisation.}
If the conversion percentage is small, we can approximate: assume a fixed
base Toughness (from class stats) and compute the trait bonus from that base
alone, ignoring the gear contribution to the source stat.  This gives a
linear approximation that can be used in the LP relaxation.

\paragraph{Approach 2: Iterative solving.}
Solve the problem without trait conversions.  Use the resulting stat totals
to compute trait bonuses.  Add these bonuses as flat adjustments and re-solve.
Repeat until the solution stabilises.  This converges for small conversion
percentages (contraction mapping argument).

\paragraph{Approach 3: Heuristic-only.}
For complex trait interactions (multiple conversions that chain), fall back
to heuristic methods (SA, GA) that don't require linearity.  Use the exact
solver on the simplified (no-trait) problem to warm-start.

\begin{warningbox}
Trait interactions can make the optimisation problem non-convex even in the
continuous relaxation.  The gap between ``solve ignoring traits'' and ``solve
with traits'' can be significant for builds that stack a single stat for
conversion (e.g., full Toughness gear with a Toughness$\to$Power trait).
Always sanity-check the solver's output against in-game stat screens.
\end{warningbox}

\begin{exercisebox}
Implement the iterative solving approach (Approach~2) for a single trait
conversion: 7\% Toughness $\to$ Power.  How many iterations until convergence
(within 1 stat point)?  Prove termination using the contraction mapping theorem.
\end{exercisebox}

\section{What-If Mode}

A player often has some gear already and wants to optimise the rest.
``What-if'' mode fixes some variables before the solver runs:

\begin{leanbox}[title={Lean 4 --- What-If Mode}]
\begin{minted}{lean}
/-- Fix specific slots to specific combos before solving. -/
def applyWhatIf (domains : GearSlot → Finset StatCombo)
    (fixed : List (GearSlot × StatCombo))
    : GearSlot → Finset StatCombo :=
  fun slot =>
    match fixed.find? (fun (s, _) => s == slot) with
    | some (_, combo) => {combo}  -- singleton domain
    | none => domains slot

-- Example: "I have Legendary Berserker armour"
def legendaryBerserkerArmour : List (GearSlot × StatCombo) :=
  [(.headgear, berserkers), (.shoulders, berserkers),
   (.coat, berserkers), (.gloves, berserkers),
   (.leggings, berserkers), (.boots, berserkers)]
\end{minted}
\end{leanbox}

This is trivially supported by the solver architecture: fixing a slot to a
singleton domain before propagation means propagation will immediately
exploit the fixed information.  The rune propagator, for instance, will
see that all armour slots have the same rune (whatever rune Berserker armour
uses) and skip rune selection entirely.

The search space shrinks dramatically: if 6 of 16 slots are fixed, the
remaining space is $40^{10} \approx 10^{16}$ instead of $40^{16} \approx
10^{25}$---and propagation typically prunes this much further.

\begin{exercisebox}
Implement what-if mode.  Solve the following scenario: a Warrior with full
Berserker armour (fixed) wants to optimise trinkets and weapons for
Power $\geq 3000$, Precision $\geq 2200$, Ferocity $\geq 1500$.  How many
B\&B nodes does the solver explore compared to the unconstrained problem?
\end{exercisebox}

\section{Pareto Frontier Mode}

Use the scalarisation sweep from Chapter~\ref{ch:pareto} with the full
solver.  For each weight vector, solve the scalarised problem to
optimality.  Output: the set of non-dominated builds across all weight
vectors---an approximation of the Pareto frontier.

\begin{leanbox}[title={Lean 4 --- Pareto Frontier Mode}]
\begin{minted}{lean}
/-- Generate the Pareto frontier by sweeping scalarisation weights. -/
def paretoFrontierMode (problem : GearProblem)
    (statA statB : StatType) (nPoints : Nat := 20)
    : List (Build × Nat × Nat) :=
  (List.range (nPoints + 1)).filterMap fun i =>
    let w := (i.toFloat) / nPoints.toFloat  -- weight for statA miss
    let weightedProblem := { problem with
      wMissPerStat := fun s =>
        if s == statA then (w * 100).toNat
        else if s == statB then ((1.0 - w) * 100).toNat
        else 50  -- default weight for other stats }
    match solve weightedProblem with
    | .optimal build pen =>
      some (build, statTotal build statA, statTotal build statB)
    | .infeasible => none

/-- Filter to non-dominated solutions. -/
def filterDominated (results : List (Build × Nat × Nat))
    : List (Build × Nat × Nat) :=
  results.filter fun (_, a1, b1) =>
    ¬ results.any fun (_, a2, b2) =>
      (a2 ≥ a1 ∧ b2 ≥ b1) ∧ (a2 > a1 ∨ b2 > b1)
\end{minted}
\end{leanbox}

Each point on the frontier represents a genuinely different gear philosophy.
On the Power/Vitality frontier, one end is full Berserker (maximum damage,
minimum health); the other is full Soldier or Minstrel (maximum health,
minimum damage).  The interior points are the interesting ones---mixed
builds that trade damage for survivability at different rates.

\begin{intuitionbox}
The Pareto frontier mode is expensive: it runs the full solver $n$ times.
But most of the work is shared: the propagation at the root node is identical
for all weight vectors (it depends on stat targets, not weights), and the
warm-start from SA can be reused across adjacent weight vectors.
\end{intuitionbox}

\section{Real-World Analogues}

The solver architecture---propagation $\to$ B\&B $\to$ LP bounding
$\to$ heuristic warm-start---is not specific to GW2.  It is the core
design of industrial solvers like Gurobi, CPLEX, and Google OR-Tools.
Problems with the same structure include:

\begin{itemize}[nosep]
  \item \textbf{Portfolio optimisation:} choose discrete asset lots to
    meet return targets at minimum risk.
  \item \textbf{Cloud provisioning:} select VM types to meet resource
    requirements at minimum cost.
  \item \textbf{Sports team composition:} select players within a salary
    cap to maximise expected performance.
  \item \textbf{Hardware configuration:} choose components (CPU, RAM,
    storage) to meet workload requirements.
\end{itemize}

In each case, the decision variables are discrete, constraints couple them,
and the objective involves meeting targets with asymmetric penalties.

\begin{intuitionbox}
You haven't just built a gear optimiser.  You've built a miniature
constraint programming / mixed-integer programming solver, verified in
Lean~4.  The techniques generalise to any problem with the same structure.
\end{intuitionbox}

% ════════════════════════════════════════════════════════════════════════════
%  APPENDICES
% ════════════════════════════════════════════════════════════════════════════
\appendix

\chapter{WvW Stat Combo Reference}
\label{app:combos}

The 40 WvW-legal stat combinations at Ascended rarity.

\section{Triple-Attribute (1 Major, 2 Minor)}

\begin{center}
\small
\begin{longtable}{llll}
\toprule
\textbf{Name} & \textbf{Major} & \textbf{Minor 1} & \textbf{Minor 2} \\
\midrule
\endhead
Berserker's & Power & Precision & Ferocity \\
Zealot's & Power & Precision & Healing Power \\
Soldier's & Power & Toughness & Vitality \\
Valkyrie & Power & Vitality & Ferocity \\
Harrier's & Power & Healing Power & Concentration \\
Captain's & Precision & Power & Toughness \\
Rampager's & Precision & Power & Condition Damage \\
Assassin's & Precision & Power & Ferocity \\
Knight's & Toughness & Power & Precision \\
Cavalier's & Toughness & Power & Ferocity \\
Nomad's & Toughness & Vitality & Healing Power \\
Settler's & Toughness & Condition Damage & Healing Power \\
Giver's & Toughness & Healing Power & Concentration \\
Sentinel's & Vitality & Power & Toughness \\
Shaman's & Vitality & Condition Damage & Healing Power \\
Sinister & Condition Damage & Power & Precision \\
Carrion & Condition Damage & Power & Vitality \\
Rabid & Condition Damage & Precision & Toughness \\
Dire & Condition Damage & Toughness & Vitality \\
Apostate's & Condition Damage & Toughness & Healing Power \\
Bringer's & Expertise & Precision & Vitality \\
Cleric's & Healing Power & Power & Toughness \\
Magi's & Healing Power & Precision & Vitality \\
Apothecary's & Healing Power & Toughness & Condition Damage \\
\bottomrule
\end{longtable}
\end{center}

\section{Quadruple-Attribute (2 Major, 2 Minor)}

\begin{center}
\small
\begin{longtable}{lllll}
\toprule
\textbf{Name} & \textbf{Major 1} & \textbf{Major 2} & \textbf{Minor 1} & \textbf{Minor 2} \\
\midrule
\endhead
Commander's & Power & Precision & Toughness & Concentration \\
Demolisher & Power & Precision & Toughness & Ferocity \\
Marauder & Power & Precision & Vitality & Ferocity \\
Vigilant & Power & Toughness & Concentration & Expertise \\
Crusader & Power & Toughness & Ferocity & Healing Power \\
Wanderer's & Power & Vitality & Toughness & Concentration \\
Diviner's & Power & Concentration & Precision & Ferocity \\
Dragon's & Power & Ferocity & Precision & Vitality \\
Viper's & Power & Condition Damage & Precision & Expertise \\
Grieving & Power & Condition Damage & Precision & Ferocity \\
Marshal's & Power & Healing Power & Precision & Condition Damage \\
Seraph & Precision & Condition Damage & Concentration & Healing Power \\
Trailblazer's & Toughness & Condition Damage & Vitality & Expertise \\
Minstrel's & Toughness & Healing Power & Vitality & Concentration \\
Ritualist's & Vitality & Condition Damage & Concentration & Expertise \\
Plaguedoctor's & Vitality & Condition Damage & Healing Power & Concentration \\
\bottomrule
\end{longtable}
\end{center}

\section{Celestial (WvW Version)}

All seven stats receive equal coefficients:
Power, Precision, Ferocity, Vitality, Toughness, Healing Power, Condition Damage.

\medskip\noindent
\textbf{Note:} The WvW version of Celestial excludes Expertise and
Concentration, which are included in the PvE version.

\chapter{Cross-Reference Map}
\label{app:crossref}

\begin{center}
\small
\begin{longtable}{p{5.5cm}p{7cm}}
\toprule
\textbf{Lattice Book Concept} & \textbf{This Book Application} \\
\midrule
\endhead
Cell with join-semilattice values &
  Gear slot domains (\texttt{Finset StatCombo}, reverse inclusion) \\
\addlinespace
Propagator (monotone function) &
  Rune, LinearGeq, LinearLeq, infusion-budget propagators \\
\addlinespace
Knaster--Tarski fixed point &
  Propagation termination + correctness \\
\addlinespace
Belnap's 4-valued lattice &
  3-valued SAT assignment lattice \\
\addlinespace
Product lattice &
  Combined gear-slot domain lattice \\
\addlinespace
Detecting $\top$ (contradiction) &
  Domain goes empty $\to$ prune this branch \\
\addlinespace
Monotone convergence &
  Each B\&B node runs propagation to fixpoint \\
\addlinespace
\texttt{Cell.update} (join new info) &
  Propagating constraints narrows domains \\
\addlinespace
Lattice height = termination bound &
  Total domain size bounds propagation steps \\
\addlinespace
Set-of-possible-values lattice &
  Exactly the domain-reduction lattice for gear slots \\
\addlinespace
Partial order &
  Pareto dominance, subproblem refinement, vertex ordering \\
\bottomrule
\end{longtable}
\end{center}

% ════════════════════════════════════════════════════════════════════════════
\end{document}
